\documentclass[11pt]{article}

% ---------------------------------------------------------------------------
% Portable preamble: compiles with pdfLaTeX or XeLaTeX, no project macros.
% ---------------------------------------------------------------------------
\usepackage{amsmath,amsthm,amssymb}
\usepackage[margin=1in]{geometry}
\usepackage[colorlinks=true,linkcolor=blue,urlcolor=blue,citecolor=blue]{hyperref}
\usepackage{mathtools}
\usepackage{enumitem}
\usepackage{booktabs}
\usepackage{longtable}
\usepackage[round,authoryear]{natbib}

\theoremstyle{plain}
\newtheorem{theorem}{Theorem}[section]
\newtheorem{lemma}[theorem]{Lemma}
\newtheorem{proposition}[theorem]{Proposition}
\newtheorem{corollary}[theorem]{Corollary}
\theoremstyle{definition}
\newtheorem{assumption}{Assumption}
\newtheorem{definition}[theorem]{Definition}
\theoremstyle{remark}
\newtheorem{remark}[theorem]{Remark}
\newtheorem{example}[theorem]{Example}
\newtheorem{convention}[theorem]{Convention}

\renewcommand{\theassumption}{A\arabic{assumption}}

\newcommand{\E}{\mathbb{E}}
\newcommand{\PP}{\mathbb{P}}
\newcommand{\Pn}{\widehat{\PP}_n}
\newcommand{\R}{\mathbb{R}}
\newcommand{\cX}{\mathcal{X}}
\newcommand{\cS}{\mathcal{S}}
\newcommand{\cY}{\mathcal{Y}}
\newcommand{\cO}{\mathcal{O}}
\newcommand{\cM}{\mathcal{M}}
\newcommand{\cQ}{\mathcal{Q}}
\newcommand{\cB}{\mathcal{B}}
\newcommand{\cK}{\mathcal{K}}
\newcommand{\Q}{\mathcal{Q}}
\newcommand{\Qhat}{\widehat{\mu}_M}
\newcommand{\ez}{e_0}
\newcommand{\ehat}{\widehat{e}}
\newcommand{\muhat}{\widehat{\mu}}
\newcommand{\what}{\widehat{w}}
\newcommand{\that}{\widehat{\tau}}
\newcommand{\Thetahat}{\widehat{\Theta}}
\newcommand{\pihat}{\widehat{\pi}}
\newcommand{\Sighat}{\widehat{\Sigma}}
\newcommand{\Ltwo}[1]{\lVert #1 \rVert_{2}}
\newcommand{\Lw}[1]{\lVert #1 \rVert_{2,w}}
\newcommand{\Linf}[1]{\lVert #1 \rVert_{\infty}}
\newcommand{\IF}{\mathrm{IF}}
\newcommand{\Cov}{\mathrm{Cov}}
\newcommand{\Corr}{\mathrm{cor}}
\DeclareMathOperator*{\argmin}{arg\,min}
\DeclareMathOperator*{\argmax}{arg\,max}
\DeclareMathOperator{\Var}{Var}
\DeclareMathOperator{\tr}{tr}
\DeclareMathOperator{\plim}{plim}
\DeclarePairedDelimiter{\abs}{\lvert}{\rvert}
\DeclarePairedDelimiter{\card}{\lvert}{\rvert}
\DeclarePairedDelimiter{\norm}{\lVert}{\rVert}

\title{Evaluating Surrogate Transportability via Local Geometric Analysis\\[2pt]
\large Theory}
\author{}
\date{}

\begin{document}
\maketitle

\begin{abstract}
We develop inference for the correlation of treatment effects on a surrogate marker
and a clinical outcome across a distribution of hypothetical future studies, drawn
uniformly from a total-variation ball of radius $\lambda$ around the observed study
$\PP_0$. Our main result is a $\sqrt n$-consistent, asymptotically normal estimator of
this correlation on a finite (or discretized) covariate space, built from
importance-weighted or cross-fitted augmented-inverse-probability-weighted (AIPW)
treatment-effect estimates in each of $M$ hit-and-run Monte Carlo draws from the ball,
with a two-stage influence function combining $\sqrt n$ estimation error and
$\sqrt M$ Monte Carlo error. The randomized-trial estimator is recovered as an exact
corollary of the AIPW theorem. We further give a closed-form characterization of the
estimand when the ball's probability-simplex boundary constraints do not bind: in this
interior regime the correlation reduces exactly, independently of $\lambda$ and of
$\PP_0$'s composition, to the ordinary Pearson correlation of the conditional average
treatment effects across covariate strata --- the same object a stratum-level
meta-analysis would report. We extend the finite-support theory to a general smooth
functional of the across-study treatment-effect distribution on a possibly continuous
covariate space, showing that every load-bearing moment is a linear or bilinear
functional of the two conditional average treatment effect (CATE) surfaces, derive the
efficient influence function of such a bilinear functional, and give the smoothness
condition under which $\sqrt n$-consistent estimation remains possible. The
finite-support theorem is recovered as the discrete special case of this general
theory. We state throughout, and discuss the consequences of, the causal scope
restriction that future studies differ from $\PP_0$ only in covariate composition, and
we condition on the sampled geometry, deferring a fully unconditional treatment to
future work.
\end{abstract}

\tableofcontents

% ===========================================================================
\section{Scope and organisation}\label{sec:scope}
% ===========================================================================

This document is the theoretical companion to the manuscript
\texttt{inst/paper/main.tex} (``Evaluating Surrogate Transportability via Local
Geometric Analysis'').

\paragraph{What this document covers.} Two co-equal main results, built on the same
influence-function machinery.

\begin{itemize}
  \item \textbf{The finite-support theorem} (Section~\ref{sec:layer1}): $X$ takes
  finitely many values (e.g., after discretizing a continuous covariate). This is the
  case implemented in the \texttt{surrogateTransportability} package
  (\texttt{tv\_ball\_correlation\_IF\_adaptive()}) and validated by the four-DGP
  simulation study in the main text. We state one unified theorem for the
  cross-fitted AIPW estimator (Theorem~\ref{thm:main}), with the randomized-trial,
  importance-weighting estimator recovered as an exact corollary
  (Corollary~\ref{cor:rct}) rather than as a separately re-derived special case. A
  further result (Proposition~\ref{prop:interior}) gives the closed-form,
  $\lambda$-free reduction of the estimand in the interior of the probability simplex.
  \item \textbf{The general theorem} (Section~\ref{sec:layer2}): $X$ is allowed to be
  continuous, and the target functional $\Theta$ is allowed to be any smooth function
  of finitely many moments of the across-study effect distribution (correlation,
  $R^2$, mean squared prediction error, and concordance are instances). This section
  is not implemented in the package and not exercised by the simulation study; it is
  included because it strictly generalizes Section~\ref{sec:layer1}, and
  Section~\ref{sec:recovery} verifies the finite-support theorem is recovered as its
  discrete special case.
\end{itemize}

\paragraph{Standing convention.} Both sections \emph{condition on the sampled
geometry}: the $M$ Monte Carlo draws $\{\Q_m\}$ (equivalently, the induced
reweighting-covariance kernel) are treated as fixed inputs to the inference problem
for the data $O_1,\dots,O_n$, and the two sources of randomness --- estimation error in
$n$ and Monte Carlo error in $M$ --- are combined explicitly rather than assumed away.
We further condition on using a $\sqrt n$-consistent estimate $\Pn$ of $\PP_0$ to
\emph{center} the sampled ball; a fully unconditional theory that additionally
differentiates the geometry's own dependence on $\PP_0$ is deferred to future work
(Remark~\ref{rem:geometry-scope}). Every result inherits this convention; we do not
repeat it in each statement.

\paragraph{What is proved in full vs.\ what is scoped out.} Every theorem, proposition,
and corollary in the main text of this document has a complete proof in
Appendix~\ref{app:proofs}. Section~\ref{sec:layer2}'s Corollary~\ref{cor:onestep-match}
is the one point in this document where we assert a structural fact --- the general
one-step bilinear correction and the package's finite-support ad hoc debiasing
correction are the same object up to normalization --- without closing the exact
per-cell normalization bookkeeping required for a numerically exact match;
Remark~\ref{rem:onestep-scope} states this limitation precisely.
Appendix~\ref{app:counterex} collects four counterexamples delimiting the theory's
scope.

% ===========================================================================
\section{Setup, notation, and standing assumptions}\label{sec:setup}
% ===========================================================================

\subsection{Data, treatment effects, and the estimand}

We observe $n$ i.i.d.\ copies $O_1,\dots,O_n$ of $O=(X,A,S,Y)\sim\PP_0$: $A\in\{0,1\}$
is a binary treatment, $S$ is a surrogate marker, $Y$ is the clinical outcome of
interest, and $X\in\cX$ are pre-treatment covariates. For any probability measure $\PP$
on the same sample space, define the treatment effects on the surrogate and the
outcome,
\begin{equation}\label{eq:delta-def}
    \Delta_S(\PP) = \E_\PP\{S(1)-S(0)\}, \qquad \Delta_Y(\PP) = \E_\PP\{Y(1)-Y(0)\}.
\end{equation}
Under randomization or unconfoundedness (Assumption~\ref{ass:causal} below), these are
identified from the observables under $\PP$.

A \emph{future study} is a probability measure $\Q$ on the same sample space. We model
$\Q$ as a draw from a probability measure $\mu$ (a distribution over study
distributions), supported on a \emph{local geometry}
\begin{equation}\label{eq:geometry}
    U(\PP_0,\lambda;d) = \{\Q\in\cM_1(\cO): d(\Q,\PP_0)\le\lambda\}
\end{equation}
around $\PP_0$, where $d$ is a distance metric between probability measures and
$\lambda\ge0$ is a tolerance. Our primary instantiation takes $d$ to be total variation,
$\mu=\mathrm{Uniform}\{U(\PP_0,\lambda;d_{\mathrm{TV}})\}$, and the estimand to be the
correlation of the two treatment effects across studies,
\begin{equation}\label{eq:estimand}
    \Theta(\PP_0,\lambda) = \Corr_{\Q\sim\mu}\{\Delta_S(\Q),\Delta_Y(\Q)\}
    = \frac{\Cov_\mu\{\Delta_S(\Q),\Delta_Y(\Q)\}}
           {\sqrt{\Var_\mu\{\Delta_S(\Q)\}\,\Var_\mu\{\Delta_Y(\Q)\}}}.
\end{equation}
$\Theta\approx1$ means the surrogate transports well across the population
compositions $\mu$ deems plausible; $\Theta\approx0$ means it does not.

\subsection{Absolute continuity and the signed-measure parametrization}

We restrict attention to $\Q\ll\PP_0$: future studies reweight the observed support of
$\PP_0$ rather than extrapolating to unobserved covariate values. It is convenient to
parametrize $\Q$ directly as a perturbation of $\PP_0$. Write
\begin{equation}\label{eq:signed-measure}
    \Q = \PP_0 + V,
\end{equation}
where $V$ is a finite signed measure with $V(\cO)=0$ (so $\Q$ is a probability measure)
and $\PP_0+V\ge0$ pointwise (so $\Q$ is nonnegative). For the total-variation ball,
$d_{\mathrm{TV}}(\Q,\PP_0)=\tfrac12\|V\|_1\le\lambda$, i.e.\ $\|V\|_1\le2\lambda$. The
density ratio (importance weight) is $w(x)=d\Q/d\PP_0(x) = 1+v(x)/p_0(x)$ where $v$ is
the density of $V$ against a base measure containing $\PP_0$; when $\cX$ is finite this
is simply $w(k)=q_k/p_0(k)$ for cell probabilities $q_k=\PP_0(k)+V(k)$.

This parametrization makes explicit that $U(\PP_0,\lambda;d_{\mathrm{TV}})$, restricted
to the finite-support case, is the intersection of three convex constraint sets: the
zero-sum hyperplane $\{V:\mathbf{1}^\top V=0\}$, the $\ell_1$-ball
$\{V:\|V\|_1\le2\lambda\}$, and the simplex-nonnegativity region
$\{V:\PP_0+V\ge0\}$. Only the first two are permutation-symmetric in the cells; the
third depends on $\PP_0$'s composition and is what makes the general theory
$\lambda$- and $\PP_0$-dependent (Section~\ref{sec:layer1}), while its absence in a
restricted range of $\lambda$ is what makes the closed-form reduction of
Proposition~\ref{prop:interior} possible.

\subsection{Conditional average treatment effects}

Define the conditional average treatment effects (CATEs)
\begin{equation}\label{eq:cate}
    \tau_S(x) = \E_{\PP_0}\{S(1)-S(0)\mid X=x\}, \qquad
    \tau_Y(x) = \E_{\PP_0}\{Y(1)-Y(0)\mid X=x\}.
\end{equation}

\begin{assumption}[A0: Conditional-effect transportability given $X$]\label{ass:A0}
Future studies $\Q\in\mathrm{supp}(\mu)$ differ from $\PP_0$ only in the marginal law of
$X$: the reweighting $w=d\Q/d\PP_0$ is a function of $X$ alone, and $\tau_S,\tau_Y$ are
the same under $\Q$ and $\PP_0$. Equivalently, all cross-study variation in treatment
effects is compositional in the measured covariates.
\end{assumption}

Under Assumption~\ref{ass:A0}, study-level effects are exactly the linear functionals
\begin{equation}\label{eq:delta-linear}
    \Delta_a(\Q) = \E_{\PP_0}\{w(X)\,\tau_a(X)\}, \qquad a\in\{S,Y\}.
\end{equation}

\begin{remark}[Definitional vs.\ interpretive reading of A0]\label{rem:A0-reading}
Under a \emph{definitional} reading, $\Theta(\PP_0,\lambda)$ simply \emph{is} the
transportability functional over the class of $X$-compositional future studies, and
Assumption~\ref{ass:A0} is a scope statement rather than a restriction. Under an
\emph{interpretive} reading --- treating $\Theta$ as answering ``does the surrogate
transport to arbitrary real future studies?'' --- Assumption~\ref{ass:A0} is a
substantive causal assumption (effect modification by $X$ is the only source of
cross-study heterogeneity), it is not testable from a single study, and it cannot be
relaxed to a point estimate without additional data. See
Section~\ref{sec:discussion}.
\end{remark}

\subsection{The reweighting-covariance kernel}

A geometry induces a mean reweighting function and a \emph{reweighting-covariance
kernel},
\begin{equation}\label{eq:kernel-def}
    \bar a(x) = \E_\mu[w(x)], \qquad C(x,x') = \Cov_\mu\{w(x),w(x')\}.
\end{equation}
Taking $\mu$-moments of \eqref{eq:delta-linear} (Fubini, justified by bounded weights
and $\E_{\PP_0}[\tau_a^2]<\infty$) reduces every load-bearing moment of the estimand to
a linear or bilinear functional of the CATE pair $(\tau_S,\tau_Y)$:
\begin{align}
    \E_\mu[\Delta_a(\Q)] &= \int \bar a(x)\,\tau_a(x)\,d\PP_0(x),
    \tag{linear}\label{eq:reduction-linear}\\
    \Cov_\mu\{\Delta_a(\Q),\Delta_b(\Q)\} &= \iint C(x,x')\,\tau_a(x)\,\tau_b(x')\,
        d\PP_0(x)\,d\PP_0(x') \;=:\; \Theta_{ab}. \tag{bilinear}\label{eq:reduction-bilinear}
\end{align}
On a finite covariate space $\cX=\{1,\dots,K\}$ with cell probabilities $p_0(k)$, write
$\tau_a=(\tau_a(1),\dots,\tau_a(K))^\top$; the kernel collapses to a $K\times K$ matrix
\begin{equation}\label{eq:sigma-def}
    \Sigma_{kk'} = \Cov_\mu(q_k,q_{k'}) = C(k,k')\,p_0(k)\,p_0(k'),
\end{equation}
and
\begin{equation}\label{eq:discrete-quadform}
    \Theta_{ab} = \tau_a^\top\Sigma\,\tau_b.
\end{equation}
This is the object Sections~\ref{sec:layer1}--\ref{sec:interior} work with directly;
Section~\ref{sec:layer2} works with the general kernel $C$.

\subsection{Standing assumptions}

\begin{assumption}[Sampling]\label{ass:iid}
$O_1,\dots,O_n$ are i.i.d.\ draws from $\PP_0$.
\end{assumption}

\begin{assumption}[Identification and positivity]\label{ass:causal}
For $a\in\{0,1\}$: $Y(a),S(a)\perp A\mid X$ under $\PP_0$ (randomization or
unconfoundedness), and $\PP_0(A=a\mid X)>0$ almost surely. In the observational case,
the propensity score satisfies $e(x)=\PP_0(A=1\mid X=x)\in[c,1-c]$ for some $c>0$
(strong positivity).
\end{assumption}

\begin{assumption}[Moments and bounded weights]\label{ass:moments}
$\E_{\PP_0}[S^2+Y^2]<\infty$, and there exists $C_w<\infty$ such that
$\sup_{\Q\in\mathrm{supp}(\mu)}\sup_{x\in\mathrm{supp}(\PP_0)} w(x)\le C_w$.
\end{assumption}

\begin{assumption}[Non-degeneracy]\label{ass:nondegeneracy}
There exists $c_0>0$ such that $\Var_\mu\{\Delta_S(\Q)\}\ge c_0$ and
$\Var_\mu\{\Delta_Y(\Q)\}\ge c_0$.
\end{assumption}

\begin{assumption}[Nuisance rates, cross-fitted case]\label{ass:nuisance}
Propensity and outcome-regression estimators $\ehat,\hat\mu_a^S,\hat\mu_a^Y$ ($a\in\{0,1\}$)
are obtained by $K$-fold cross-fitting and satisfy
$\|\ehat-e\|_{L^2(\PP_0)}=o_P(n^{-1/4})$,
$\|\hat\mu_a-\mu_a\|_{L^2(\PP_0)}=o_P(n^{-1/4})$, hence the product rate
$\|\ehat-e\|_{L^2}\cdot\|\hat\mu_a-\mu_a\|_{L^2}=o_P(n^{-1/2})$.
\end{assumption}

\begin{assumption}[Geometry regularity and MCMC ergodicity]\label{ass:geometry}
$\mathrm{supp}(\mu)$ is compact, and the hit-and-run sampler generating
$\{\Q_1,\dots,\Q_M\}$ is geometrically ergodic with stationary distribution $\mu$.
\end{assumption}

\begin{convention}[Conditioning on the geometry]\label{conv:geometry}
Throughout, inference is conditional on the sampled geometry $\{\Q_1,\dots,\Q_M\}$
(equivalently on the empirical kernel it induces). This is the convention flagged in
Section~\ref{sec:scope}; Remark~\ref{rem:geometry-scope} states what is deferred.
\end{convention}

\begin{remark}[Scope: the geometry is centered at an estimate of $\PP_0$]\label{rem:geometry-scope}
In practice $\{\Q_m\}$ are drawn from $U(\Pn,\lambda;d_{\mathrm{TV}})$, a ball centered at
the empirical distribution $\Pn$, not at $\PP_0$. Two consequences, both handled by
Convention~\ref{conv:geometry} rather than resolved: (i) \emph{estimation, conditional
on a drawn $\Q_m$, is unaffected.} For the finite-support estimators of
Section~\ref{sec:layer1}, the empirical cell probability entering the importance
weight cancels exactly against the per-cell sample-size normalization in the
treatment-effect estimator (Lemma~\ref{lem:cancellation}), so $\hat\Delta_a(\Q_m)$ is a
consistent, asymptotically linear estimator of the \emph{true}
$\Delta_a(\Q_m)=\sum_k q_{m,k}\tau_a(k)$ using the true CATEs, for any fixed $\Q_m$,
regardless of whether $\Q_m$ was itself drawn from a ball centered at $\PP_0$ or
$\Pn$. (ii) \emph{which $\Q_m$'s get drawn is data-dependent.} Because
$U(\Pn,\lambda)\to U(\PP_0,\lambda)$ as $\Pn\to\PP_0$ (the ball's boundary is Lipschitz
in its center for the total-variation metric), $\hat\mu_M\to\mu$ picks up an additional
$O_P(n^{-1/2})$ recentering error beyond the $O_P(M^{-1/2})$ Monte Carlo error of
Lemma~\ref{lem:mcmc-error}. A fully unconditional theory would differentiate through
this dependence (as it would through the kernel $C$'s dependence on $\PP_0$ in
Section~\ref{sec:layer2}, Remark~\ref{rem:kernel-dependence}); we do not pursue this
here, and the conditional statements of Theorems~\ref{thm:main} and~\ref{thm:general}
are unaffected by it.
\end{remark}

% ===========================================================================
\section{The interior-regime reduction: a closed form}\label{sec:interior}
% ===========================================================================

Before developing the general estimation theory, we characterize what the estimand
\emph{is} in the simplest case, where the local geometry's simplex-boundary
constraints do not bind. This gives an immediate interpretive anchor: in this regime,
$\Theta$ is exactly the stratum-level meta-analytic correlation of treatment effects,
computed within a single study rather than across several.

\begin{definition}[Interior regime]\label{def:interior}
The local geometry $U(\PP_0,\lambda;d_{\mathrm{TV}})$ on a finite covariate space
$\cX=\{1,\dots,K\}$ is in the \emph{interior regime} if the nonnegativity constraint
$\PP_0+V\ge0$ never binds on the feasible set, i.e.\ if
$\{V:\mathbf 1^\top V=0,\ \|V\|_1\le2\lambda\} = \{V:\mathbf 1^\top V=0,\ \|V\|_1\le2\lambda,\ \PP_0+V\ge0\}$.
A sharp sufficient condition is $\lambda\le\min_k p_0(k)$: for any $V$ with
$\mathbf1^\top V=0$ and $\|V\|_1\le2\lambda$, the triangle inequality
$\|V\|_1\ge|V_k|+|\sum_{j\ne k}V_j|=2|V_k|$ (using the zero-sum constraint) forces
$|V_k|\le\lambda$ for every $k$, so $q_k=p_0(k)+V_k\ge p_0(k)-\lambda\ge0$ whenever
$\lambda\le\min_kp_0(k)$; this bound is attained (not merely sufficient) when the
$\ell_1$ budget is split between cell $k$ and a single other cell.
\end{definition}

\begin{lemma}[The interior regime requires atomic support, and the sufficient
condition is sharp]\label{lem:interior-atomic}
Fix $\lambda\in(0,1]$ (the total-variation distance is bounded by $1$, so
$\lambda>1$ is vacuous; also exclude the degenerate $K=1$ case, where $V\equiv0$ is
forced by $V\ll\PP_0$ and $V(\cX)=0$, and the interior regime holds trivially at
every $\lambda$). Definition~\ref{def:interior} extends verbatim to a general
covariate law $\PP_0$ on any measurable $\cX$, not only a finite one, once the
signed-measure parametrization of \eqref{eq:signed-measure} is read literally (with
$V\ll\PP_0$ throughout, the ambient constraint of \S2.2): $U(\PP_0,\lambda;d_{\mathrm{TV}})$
is in the interior regime at $\lambda$ if $\{V:V(\cX)=0,\|V\|_1\le2\lambda\}\subseteq\{V:\PP_0+V\ge0\}$.
This holds at $\lambda$ if and only if there is no measurable $A\subseteq\cX$ with
$0<\PP_0(A)<\lambda$; equivalently (Step 0 below), if and only if $\PP_0$ is purely
atomic with finitely many atoms $\{1,\dots,K\}$, $K\le\lfloor1/\lambda\rfloor$, and
$\min_k p_0(k)\ge\lambda$. In particular, Definition~\ref{def:interior}'s sufficient
condition $\lambda\le\min_kp_0(k)$ is not merely sufficient but \emph{sharp}.
\end{lemma}

\begin{proof}
\emph{Step 0 (the two characterizations coincide).} If $\PP_0$ has a non-atomic
component of mass $c>0$, that component's range $\{\PP_0(B):B\subseteq\text{component}\}$
is the full interval $[0,c]$ (a standard consequence of countable additivity for
non-atomic measures, e.g.\ via repeated halving --- the Sierpi\'nski/Lyapunov range
theorem), so it contains a set $A$ with $0<\PP_0(A)<\min(\lambda,c)\le\lambda$. If
$\PP_0$ has infinitely many atoms, their masses sum to at most $1$ and hence tend to
$0$, so some atom has mass $<\lambda$ and is itself such an $A$. Conversely, if
$\PP_0$ is purely atomic with finitely many atoms all of mass $\ge\lambda$, no such
$A$ exists (any measurable set is a union of atoms, each contributing $0$ or
$\ge\lambda$ to its total mass, so no partial sum lands strictly between $0$ and
$\lambda$); finiteness of the atom count is then forced by $\sum_kp_0(k)=1$ and
$p_0(k)\ge\lambda>0$, giving $K\le\lfloor1/\lambda\rfloor$. So ``no such $A$ exists''
$\iff$ ``$\PP_0$ purely atomic, finitely many atoms, $\min_kp_0(k)\ge\lambda$.''

($\Leftarrow$) If every measurable $A$ has $\PP_0(A)=0$ or $\PP_0(A)\ge\lambda$
(equivalently, by Step 0, $\PP_0$ is purely atomic with $\min_kp_0(k)\ge\lambda$),
then for any feasible $V$ ($V(\cX)=0$, $\|V\|_1\le2\lambda$) and any atom $k$, the
triangle-inequality argument already used in Definition~\ref{def:interior} gives
$|V_k|\le\lambda\le p_0(k)$, so $p_0(k)+V_k\ge0$: the interior regime holds.

($\Rightarrow$) Suppose $\exists A$ with $0<\PP_0(A)=\varepsilon<\lambda$ (so
$\varepsilon<\lambda\le1$, hence $\varepsilon<1$ and $\PP_0(A^c)=1-\varepsilon>0$).
The density perturbation $v=-\mathbf1_A+\big(\varepsilon/(1-\varepsilon)\big)\mathbf1_{A^c}$
against $\PP_0$ defines $V=v\,d\PP_0$; since $v$ is a bounded density against
$\PP_0$, $V\ll\PP_0$ throughout, respecting the ambient constraint. It satisfies
$\int v\,d\PP_0=-\varepsilon+\frac{\varepsilon}{1-\varepsilon}\cdot\PP_0(A^c)=-\varepsilon+\frac{\varepsilon}{1-\varepsilon}(1-\varepsilon)=-\varepsilon+\varepsilon=0$
(zero-sum) and drives the density to exactly $0$ on $A$. Its total-variation cost is
$\|V\|_1=\varepsilon+\varepsilon=2\varepsilon$, since the negative and positive parts
of $v\,d\PP_0$ each carry mass exactly $\varepsilon$ by construction. Because
$\varepsilon<\lambda$, $\|V\|_1=2\varepsilon<2\lambda$: this $V$ is strictly inside
the TV budget yet already drives $\PP_0+V$ to $0$ on $A$, so rescaling $v$ by a
factor $1+\delta$ for small enough $\delta>0$ (remaining feasible, since
$2\varepsilon(1+\delta)<2\lambda$ for $\delta<\lambda/\varepsilon-1>0$) drives
$\PP_0+V$ strictly negative on $A$. The interior regime therefore fails at $\lambda$.
\end{proof}

\begin{remark}[No continuous-$X$ analog of the closed form]\label{rem:no-continuous-interior}
Lemma~\ref{lem:interior-atomic} settles, in the negative, the question of whether
Proposition~\ref{prop:interior} below has an analog for continuous $X$: the interior
regime is empty for \emph{every} $\lambda>0$ whenever $\PP_0$ has a continuous
component (or infinitely many atoms), by Step 0 of the lemma's proof. The premise of
Proposition~\ref{prop:interior} --- a $\lambda>0$ at which the geometry is in the
interior regime --- therefore cannot hold outside finite atomic support, and there
is no direct continuous-$X$ closed form to derive from the same symmetry argument.
(The saturated regime of Section~\ref{sec:saturated} is a different matter: for
continuous $\PP_0$, $\sup_{\Q\ll\PP_0}d_{\mathrm{TV}}(\Q,\PP_0)=1$ still, so a
saturated regime at $\lambda\ge1$ is not ruled out by this lemma; but
Proposition~\ref{prop:saturated}'s argument does not directly generalize either,
since $\mathrm{Uniform}(\Delta_{K-1})$ has no analog on the infinite-dimensional
simplex of probability measures $\ll\PP_0$ --- there is no uniform probability
measure on an infinite-dimensional convex body, so $\mu$ itself is not well-defined
there, independently of any symmetry argument.) Section~\ref{sec:sieve-limit}
gives the honest continuous-$X$ limit instead, obtained by discretizing rather than
by relaxing Definition~\ref{def:interior} itself.
\end{remark}

\begin{proposition}[Interior-regime closed form]\label{prop:interior}
In the interior regime, for any $\lambda>0$,
\begin{equation}\label{eq:interior-reduction}
    \Theta(\PP_0,\lambda) = \Corr_K(\tau_S,\tau_Y)
    := \frac{\frac1K\sum_k(\tau_S(k)-\bar\tau_S)(\tau_Y(k)-\bar\tau_Y)}
            {\sqrt{\frac1K\sum_k(\tau_S(k)-\bar\tau_S)^2}\,\sqrt{\frac1K\sum_k(\tau_Y(k)-\bar\tau_Y)^2}},
\end{equation}
the ordinary (unweighted) Pearson correlation of the $K$ cell-level CATEs, where
$\bar\tau_a=\frac1K\sum_k\tau_a(k)$. In particular $\Theta(\PP_0,\lambda)$ does not
depend on $\lambda$ and does not depend on $\PP_0$'s composition $(p_0(k))_k$
(only on which cells have $p_0(k)>0$).
\end{proposition}

The proof (Appendix~\ref{app:proofs}, Section~\ref{app:proof-interior}) shows that in
the interior regime the feasible region $\{V:\mathbf1^\top V=0,\|V\|_1\le2\lambda\}$ is
invariant under permutations of the $K$ cells, that the uniform measure on any
permutation-invariant convex body inherits that symmetry, and that this forces the
induced covariance kernel to take the form $\Sigma=a(\lambda)(I_K-\tfrac1K\mathbf1\mathbf1^\top)$
for a scalar $a(\lambda)>0$; substituting into \eqref{eq:discrete-quadform} and forming
the ratio in \eqref{eq:estimand} cancels $a(\lambda)$ and $p_0$ entirely.

\begin{remark}[Interpretation and scope]\label{rem:interior-scope}
Proposition~\ref{prop:interior} is a genuine reduction, not an approximation, but its
regime ($\lambda\le\min_kp_0(k)$) is typically \emph{smaller} than the $\lambda$
values used for the paper's validation study (Section~\ref{sec:layer1} operates at
$\lambda=0.3$ with $\min_k p_0(k)=0.05$, well outside the interior regime; the
reported estimand there is genuinely $\lambda$- and $\PP_0$-dependent). The
proposition should be read as a limiting characterization --- what $\Theta$ converges
to as the geometry shrinks toward the interior --- that motivates the general
construction (Section~\ref{sec:layer1}) without replacing it: for $\lambda$ small
enough that no stratum can be driven to zero mass, evaluating cross-study
transportability by the local geometric method reduces to exactly what a
meta-analysis would report if each covariate stratum were an independent study; for
larger $\lambda$, the simplex boundary makes rare strata harder to reweight upward
without exhausting common strata, and the estimand reflects this asymmetry.
\end{remark}

\subsection{The interior-regime efficient influence function}\label{sec:interior-eif}

Proposition~\ref{prop:interior} identifies what $\Theta$ \emph{is} in the interior
regime; the following identifies its EIF there, specializing the general bilinear
EIF of Lemma~\ref{lem:bilinear-eif} rather than re-deriving it from scratch.

\begin{corollary}[Interior-regime efficient influence function]\label{cor:interior-eif}
Suppose $K\ge3$ and, in the interior regime with $\lambda<\min_kp_0(k)$
\emph{strictly}, the across-cell CATE spreads $s_S,s_Y>0$ and
$|\Corr_K(\tau_S,\tau_Y)|<1$ (the finite-$K$ specialization of
Assumption~\ref{ass:nondegeneracy}; $K=2$ is degenerate, since then
$\Theta=\pm1$ identically, see Remark~\ref{rem:interior-eif-scope}). Then
$\Theta(\PP_0,\lambda)=\Corr_K(\tau_S,\tau_Y)$ (Proposition~\ref{prop:interior}) is a
smooth function of the finite-dimensional parameter $(\tau_S,\tau_Y)\in\R^{2K}$,
hence pathwise differentiable, with efficient influence function
\begin{equation}\label{eq:interior-eif}
    \chi_\Theta(O) = \sum_{k=1}^K\frac{\partial\Corr_K}{\partial\tau_S(k)}(\tau_S,\tau_Y)\,\IF_{\tau_S}(O;k)
    \;+\; \sum_{k=1}^K\frac{\partial\Corr_K}{\partial\tau_Y(k)}(\tau_S,\tau_Y)\,\IF_{\tau_Y}(O;k),
\end{equation}
where $\IF_{\tau_a}(O;k)=\frac{\mathbf1\{X=k\}}{p_0(k)}\big(\mathrm{score}_a(O)-\tau_a(k)\big)$
is the per-cell CATE influence function (the pointwise, mean-zero specialization of
Lemma~\ref{lem:cate-linearity}'s $\IF_{\tau_a}$ to a point mass at $k$), and
$\mathrm{score}_a(O)$ is the cross-fitted AIPW bracket \eqref{eq:aipw-estimator}'s
per-unit expression --- a single, estimator-independent quantity (an EIF for a fixed
nonparametric model does not depend on which consistent estimator is used). On cell
$k$, since $\mu_1^a(k)-\mu_0^a(k)=\tau_a(k)$, $\mathrm{score}_a(O)-\tau_a(k)$
collapses exactly to $\frac{A(S-\mu_1^a(k))}{e(X)}-\frac{(1-A)(S-\mu_0^a(k))}{1-e(X)}$,
the \emph{stratified difference-of-means} influence function, regardless of whether
$e$ is known or estimated (the unaugmented Horvitz--Thompson score
$\frac{AS}{e(X)}-\frac{(1-A)S}{1-e(X)}-\tau_a(k)$ is a valid but strictly less
efficient influence function, not this one). \eqref{eq:interior-eif} does not depend on $a(\lambda)$: $\Sigma=a(\lambda)(I_K-\mathbf1\mathbf1^\top/K)$
enters $\Theta_{ab}=\tau_a^\top\Sigma\tau_b$ linearly and exactly once, so
$\Theta=\Theta_{SY}/\sqrt{\Theta_{SS}\Theta_{YY}}$ is homogeneous of degree $0$ in
$a(\lambda)$ (a ratio of two degree-$1$-homogeneous quantities), hence so is
$\Corr_K$'s gradient in \eqref{eq:interior-eif} --- $a(\lambda)$ is known geometry
($a(\lambda)=\lambda^2a_K$ by scaling), not an estimated nuisance, so this is an
exact cancellation, not an approximation. Equivalently, specializing
Lemma~\ref{lem:bilinear-eif}'s general EIF \eqref{eq:chi-ab} to the interior-regime
kernel $C(k,k')=a(\delta_{kk'}-1/K)/(p_0(k)p_0(k'))$ gives the intermediate,
$a$-dependent bilinear scores
\begin{equation}\label{eq:interior-chi-ab}
    \chi_{ab}(O) = a\cdot\frac{(\tau_b(X)-\bar\tau_b)\big(\mathrm{score}_a(O)-\tau_a(X)\big)}{p_0(X)}
    \;+\; (a\leftrightarrow b),
\end{equation}
and composing $\chi_{SS},\chi_{SY},\chi_{YY}$ with $\nabla_{(\Theta_{SS},\Theta_{SY},\Theta_{YY})}\Corr_K$
reproduces \eqref{eq:interior-eif} with $a$ cancelling in the composition. The same
limit also follows directly from Theorem~\ref{thm:main}: its $M$-draw influence
function \eqref{eq:psi-theta-layer1} converges, as $M\to\infty$ in the interior
regime, to \eqref{eq:interior-eif}, since $\Thetahat_n\to\Corr_K(\hat\tau_S,\hat\tau_Y)$
there, so the MCMC machinery of Section~\ref{sec:layer1} is unnecessary in this
regime.
\end{corollary}

\begin{remark}[Four subtleties in the interior-regime EIF]\label{rem:interior-eif-scope}
\textbf{Strict inequality (requires $K\ge3$).} $\lambda<\min_kp_0(k)$ is required
\emph{strictly}, and this kink is a genuine feature only for $K\ge3$: at $K=2$,
$\Theta=\pm1$ identically for every $\lambda$ and every $p_0$ (both $\Delta_a$ are
affine in the single free coordinate $q_1=1-q_2$), so $p_0\mapsto\Theta$ is constant
and trivially differentiable through the boundary. For $K\ge3$, at
$\lambda=\min_kp_0(k)$ exactly, $\PP_0\mapsto\Theta$ is continuous but not
differentiable: an active nonnegativity constraint on the argmin cell creates, on
one side of $\lambda=\min_kp_0(k)$, a $(K-2)$-dimensional feasible cross-section of
positive volume $\asymp\lambda^{K-2}$ that is absent on the other side, so moving
the boundary by $\delta$ changes that volume by $\asymp\delta\lambda^{K-2}$ (linear
in $\delta$): a genuine kink in $\Theta(p_0)$ exactly at the boundary. This is the
same combinatorial event --- the argmin cell's nonnegativity facet becoming active
--- that drives $\Theta$'s kink in $\lambda$ at \emph{fixed} $\PP_0$
(Proposition~\ref{prop:exit-interior}, Section~\ref{sec:exit-interior}): the two are
kinks along different partial directions through one geometric mechanism (here,
varying $\PP_0$ at fixed $\lambda=\min_kp_0(k)$; there, varying $\lambda$ at fixed
$\PP_0$), not two independent phenomena. This is
consistent with finite-difference slopes of $\Theta(\lambda)$ at fixed $p_0$:
statistically indistinguishable from $0$ on the interior side of
$\lambda=\min_kp_0(k)$, statistically distinguishable from $0$ on the
boundary-binding side. \textbf{Design, not test.} The interior regime should
be a design choice --- selecting $\lambda$ with margin below $\min_k\hat p_0(k)$ ---
not a data-driven hypothesis test; testing it introduces an unaccounted pre-test
affecting coverage. \textbf{Degeneracy is two-sided, not one-sided.}
\eqref{eq:interior-eif}'s gradient blows up as either across-cell CATE spread
$s_S,s_Y\to0$ (the $1/(Ks_a)$ prefactor in
$\partial\Corr_K/\partial\tau_S(k)=\frac1{Ks_S}\big[\frac{\tau_Y(k)-\bar\tau_Y}{s_Y}-\Corr_K\frac{\tau_S(k)-\bar\tau_S}{s_S}\big]$
is unbounded), addressed by the finite-$K$ specialization of
Assumption~\ref{ass:nondegeneracy} used in the corollary's hypotheses. The
\emph{opposite} boundary behaves oppositely: as $|\Corr_K|\to1$, the bracketed term
above vanishes identically for every $k$ (since $\Corr_K=\pm1$ forces the two
standardized CATE sequences to coincide up to sign), so the gradient --- and the
naive-scale asymptotic variance --- \emph{degenerates to zero}, not to infinity;
this is the familiar boundary degeneracy of a sample correlation coefficient
(classically, $\Var(r)\approx(1-\rho^2)^2/n\to0$ as $\rho\to\pm1$), and Fisher's
$z$-transform is the standard remedy for the resulting non-normal, rate-changed
sampling behavior near that boundary --- it stabilizes variance across the range of
$\rho$, not an unbounded gradient, since there is none there. \textbf{Scope: coarsened
target, not the continuum limit.} The interior condition $p_0(k)\ge\lambda$ bounds
$1/p_0(k)\le1/\lambda$ and $K\le1/\lambda$, buying $\sqrt n$-inference with no
smoothness elbow --- but only for the estimand $\Theta_K$ at the fixed, coarsened
$K$-cell partition. If the $K$ cells coarsen a richer or continuous $X$, this
describes inference for the coarsened target, not the continuum limit
$\Corr_{\PP_0}(\tau_S,\tau_Y)$ of Proposition~\ref{prop:sieve-limit}; targeting that
limit requires $K\to\infty$, and the elbow of Assumption~\ref{ass:elbow} returns via
Section~\ref{sec:rank}. This is the conceptual link between
Corollary~\ref{cor:interior-eif}, Proposition~\ref{prop:sieve-limit}, and
Section~\ref{sec:rank}: free $\sqrt n$-inference at fixed, finite $K$ is not free
$\sqrt n$-inference for the continuum target.
\end{remark}

\subsection{The saturated regime: the closed form recurs at the other extreme of
$\lambda$}\label{sec:saturated}

The interior regime is one extreme of $\lambda$ (small enough that no stratum can be
driven to zero mass); the closed form of Proposition~\ref{prop:interior} recurs,
by an independent argument, at the \emph{other} extreme, where $\lambda$ is large
enough that the local geometry covers the entire probability simplex. The
mechanism generating this recurrence does not in fact depend on which distance $d$
defined the geometry --- only on the fact that, at this radius, the geometry
\emph{equals} the whole simplex --- so we state it once, generally, before
specializing to total variation.

\begin{lemma}[Saturation is distance-free]\label{lem:saturation-general}
Fix any distance $d$ and any $\lambda$ such that $U(\PP_0,\lambda;d)=\Delta_{K-1}$
(identifying $U(\PP_0,\lambda;d)\subseteq\cM_1(\cO)$ with its image in the
covariate-marginal simplex under Assumption~\ref{ass:A0}, exactly as in the proof
below). Then $\mu_\lambda=\mathrm{Uniform}(\Delta_{K-1})=\mathrm{Dirichlet}(1,\dots,1)$,
identically for every $d$ attaining this at $\lambda$; consequently
$\Cov_{\mu_\lambda}(Q)=\frac1{K(K+1)}\big(I_K-\frac1K\mathbf1\mathbf1^\top\big)$ and
$\Theta(\PP_0,\lambda)=\Corr_K(\tau_S,\tau_Y)$.
\end{lemma}

\begin{proof}
$\mu_\lambda$ is defined as the measure uniform with respect to Lebesgue measure on
the ambient hyperplane, restricted to $U(\PP_0,\lambda;d)$: it depends on $d$ only
through this set. By hypothesis the set equals $\Delta_{K-1}$ regardless of which
$d$ produced it, so $\mu_\lambda$ is identically
$\mathrm{Uniform}(\Delta_{K-1})=\mathrm{Dir}(1,\dots,1)$ for every such $d$, which is
exchangeable under every permutation of $\{1,\dots,K\}$ (permuting the coordinates
of a symmetric Dirichlet law leaves it unchanged). Exchangeability gives
$\Cov(\Pi Q)=\Pi\Cov(Q)\Pi^\top=\Cov(Q)$ for every permutation matrix $\Pi$, so
$\Sigma=\Cov(Q)$ commutes with the full symmetric group's permutation
representation and hence, by the same commutant argument as the proof of
Proposition~\ref{prop:interior} (Appendix~\ref{app:proof-interior}), takes the form
$\Sigma=aI_K+b\mathbf1\mathbf1^\top$ for scalars $a,b$; this step uses only
exchangeability of $Q$ itself, not the centered variable $V=Q-\PP_0$ (indeed
$V\stackrel d=\Pi V$ can fail here, since the feasible set for $V$ at saturation,
$\Delta_{K-1}-\PP_0$, need not be permutation-invariant unless $\PP_0$ itself is
uniform). Since $\sum_kQ_k=1$ a.s., $\Sigma\mathbf1=\Cov(Q,\mathbf1^\top Q)=\Cov(Q,1)=0$,
which \emph{pins the form} to $b=-a/K$, giving
$\Sigma=a(I_K-\mathbf1\mathbf1^\top/K)$ for some scalar $a$; the \emph{value} of $a$
follows separately from Dirichlet$(1,\dots,1)$'s known moments,
$\mathrm{Var}(Q_k)=(K-1)/(K^2(K+1))$, $\mathrm{Cov}(Q_k,Q_{k'})=-1/(K^2(K+1))$ for
$k\ne k'$, giving $a=1/(K(K+1))>0$ (positive since $Q$ is non-degenerate for
$K\ge2$) --- though the value is not needed for the closed form, since $a$ cancels
exactly in the ratio $\Theta_{ab}=\tau_a^\top\Sigma\tau_b$,
$\Theta=\Theta_{SY}/\sqrt{\Theta_{SS}\Theta_{YY}}$, by the same
degree-zero-homogeneity cancellation as $a(\lambda)$ in
Proposition~\ref{prop:interior}.
\end{proof}

\begin{lemma}[Saturation threshold]\label{lem:saturation}
Suppose $p_0(k)>0$ for all $k=1,\dots,K$ (so every vertex of $\Delta_{K-1}$ is
reachable under the ambient constraint $\Q\ll\PP_0$). Then
$U(\PP_0,\lambda;d_{\mathrm{TV}})=\Delta_{K-1}$ (the whole probability simplex on
$K$ cells) if and only if $\lambda\ge1-\min_kp_0(k)$.
\end{lemma}

\begin{proof}
$d_{\mathrm{TV}}(\cdot,\PP_0)$ is convex on $\Delta_{K-1}$ (a norm composed with an
affine map), so its maximum over the simplex is attained at a vertex $\delta_{k^*}$;
$d_{\mathrm{TV}}(\delta_{k^*},\PP_0)=\tfrac12\big((1-p_0(k^*))+\sum_{k\ne k^*}p_0(k)\big)=1-p_0(k^*)$,
maximized over vertices by $k^*=\arg\min_kp_0(k)$. So
$\max_{\Q\in\Delta_{K-1}}d_{\mathrm{TV}}(\Q,\PP_0)=1-\min_kp_0(k)$, and the ball
covers the whole simplex exactly when $\lambda$ reaches this value.
\end{proof}

\begin{proposition}[Saturated-regime closed form]\label{prop:saturated}
If $\lambda\ge1-\min_kp_0(k)$ (Lemma~\ref{lem:saturation}), then
$\Theta(\PP_0,\lambda)=\Corr_K(\tau_S,\tau_Y)$, the same unweighted Pearson
correlation as Proposition~\ref{prop:interior}'s interior-regime reduction
\eqref{eq:interior-reduction}, again independent of $\PP_0$'s composition.
\end{proposition}

\begin{proof}
The TV instance of Lemma~\ref{lem:saturation-general}: by Lemma~\ref{lem:saturation},
$\lambda\ge1-\min_kp_0(k)$ gives $U(\PP_0,\lambda;d_{\mathrm{TV}})=\Delta_{K-1}$
(identifying $U(\PP_0,\lambda;d_{\mathrm{TV}})$ with the corresponding subset of the
covariate simplex $\Delta_{K-1}$, under Assumption~\ref{ass:A0}, exactly as
Lemma~\ref{lem:saturation-general} requires), and Lemma~\ref{lem:saturation-general}
then gives $\Theta(\PP_0,\lambda)=\Corr_K(\tau_S,\tau_Y)$ directly.
\end{proof}

\begin{remark}[The closed form is not a small-$\lambda$ artifact]\label{rem:saturated-scope}
Proposition~\ref{prop:saturated} shows the closed form of
Proposition~\ref{prop:interior} recurs at $\lambda\ge1-\min_kp_0(k)$, so it holds at
\emph{both} extremes of $\lambda$: the intermediate range
$\lambda\in(\min_kp_0(k),\,1-\min_kp_0(k))$ is the only range in which $\Theta$ need
not equal $\Corr_K(\tau_S,\tau_Y)$ and may genuinely depend on $\lambda$ and
$\PP_0$'s composition (Example~\ref{ex:binding} confirms this range is where the
paper's validated machinery does genuine work). This intermediate range is empty
exactly when $\min_kp_0(k)\ge1/2$, i.e.\ at $K=2$ with $p_0=(\tfrac12,\tfrac12)$
(where $\Theta=\pm1$ deterministically for any $\lambda$, since two numbers are
always perfectly correlated). Section~\ref{sec:discussion} revises the ``two
regimes'' framing accordingly, to three: interior, intermediate, and saturated.
\end{remark}

\subsection{The continuous-$X$ limit: an equal-mass sieve}\label{sec:sieve-limit}

Lemma~\ref{lem:interior-atomic} rules out a direct continuous-$X$ analog of
Proposition~\ref{prop:interior}. The following gives the intrinsic continuous-$X$
limit instead, obtained by refining a discretization rather than by relaxing the
interior-regime definition.

\begin{proposition}[Equal-mass dyadic sieve limit]\label{prop:sieve-limit}
Let $(\cX,\PP_0)$ be non-atomic and standard Borel, so that
$(\cX,\PP_0)\cong([0,1],\mathrm{Leb})$ under a measure isomorphism. Pull back the
dyadic intervals $\{[j2^{-m},(j+1)2^{-m})\}_{j=0}^{2^m-1}$ at level $m$ (so
$K_m=2^m$) under this isomorphism, giving a sequence of partitions that is (i)
\emph{nested} across $m$ (each level strictly refines the last), (ii) exactly equal
$\PP_0$-mass $2^{-m}$ per cell, and (iii) generates the Borel $\sigma$-field mod
$\PP_0$-null sets as $m\to\infty$ (a standard martingale-theory fact for dyadic
filtrations under a measure isomorphism to Lebesgue measure; see e.g.\
\citet{Williams1991}, Ch.\ 14). Suppose $\tau_S,\tau_Y\in L^2(\PP_0)$: this follows
from $\E[S^2],\E[Y^2]<\infty$ (Assumption~\ref{ass:moments}, stated for the observed
variables) together with strict positivity $e(x)\in[c,1-c]$
(Assumption~\ref{ass:causal}), which together give the potential-outcome moments
$\E[S(a)^2],\E[Y(a)^2]<\infty$ needed, e.g.\
$\E[S(1)^2]=\E\big[\E[S^2\mathbf1\{A=1\}\mid X]/e(X)\big]\le c^{-1}\E[S^2]<\infty$,
and $\tau_a=\mu_1^a-\mu_0^a$ is a difference of conditional expectations of such
variables. \emph{The pushforward $\PP_0^{(m)}$ of $\PP_0$ onto the level-$m$
partition} --- the discrete law on $\{1,\dots,K_m\}$ with cell probabilities
$p_0^{(m)}(k)=2^{-m}$ --- is a finite covariate law to which
Definition~\ref{def:interior} and Proposition~\ref{prop:interior} apply literally
(unlike $\PP_0$ itself, per Lemma~\ref{lem:interior-atomic}): the interior regime
holds for $\PP_0^{(m)}$ at any $\lambda\le2^{-m}=:\lambda_m$ (any choice in this
range gives the same $\Theta_m$, since $\Theta$ is $\lambda$-free throughout the
interior regime; $\lambda_m\to0$ is recorded only because the range shrinks with
$m$, not because a specific rate is needed). Setting
$\Theta_m:=\Theta(\PP_0^{(m)},\lambda_m)=\Corr_{K_m}(\tau_S^{(m)},\tau_Y^{(m)})$
with $\tau_a^{(m)}(k)=\E[\tau_a(X)\mid\text{cell }k]$ (Proposition~\ref{prop:interior}
applied to $\PP_0^{(m)}$), and assuming $\Var_{\PP_0}(\tau_S),\Var_{\PP_0}(\tau_Y)>0$
(the continuous-$X$ reading of Assumption~\ref{ass:nondegeneracy}), then
\begin{equation}\label{eq:sieve-limit}
    \Theta_m \;\longrightarrow\; \Corr_{\PP_0}\big(\tau_S(X),\tau_Y(X)\big)
    \qquad\text{as } m\to\infty,
\end{equation}
the ordinary $\PP_0$-weighted correlation of the CATE functions.
\end{proposition}

\begin{proof}
Because cells have equal $\PP_0$-mass $2^{-m}$, the unweighted average over $K_m$
cells in \eqref{eq:interior-reduction} (applied to $\PP_0^{(m)}$) coincides with the
$\PP_0$-weighted average over the same cells, so
$\Theta_m=\Corr_{\PP_0}(\tau_S^{(m)},\tau_Y^{(m)})$ exactly at each $m$ (not merely
approximately). By the martingale convergence theorem (Doob) applied to the nested,
refining dyadic filtration, $\tau_a^{(m)}\to\tau_a$ in $L^2(\PP_0)$ as $m\to\infty$.
$\Var_{\PP_0}(\tau_a^{(m)})$ is non-decreasing in $m$ (conditional-variance
decomposition: coarsening a conditioning set cannot increase the variance of the
resulting conditional expectation), so since $\Var_{\PP_0}(\tau_a)>0$ in the limit
(hypothesis), there
is a finite $m_0$ beyond which $\Var_{\PP_0}(\tau_a^{(m)})$ is bounded away from $0$
for all $m\ge m_0$ (monotone convergence to a positive limit); this is an
\emph{eventual}, not universal, non-degeneracy statement, since a coarse partition
on which $\tau_a$ happens to be constant can have $\Var_{\PP_0}(\tau_a^{(m)})=0$ at
small $m$. For $m\ge m_0$, continuity of the correlation functional under $L^2$
convergence of both arguments, given denominators bounded away from $0$, gives
\eqref{eq:sieve-limit}.
\end{proof}

\begin{remark}[Scope: canonical grid, and identification at each level]\label{rem:sieve-scope}
Two caveats delimit Proposition~\ref{prop:sieve-limit}'s scope. \emph{Grid
dependence (heuristic, backed by numerical evidence rather than a proof here).} The
limit \eqref{eq:sieve-limit} is intrinsic to the equal-$\PP_0$-mass (dyadic) grid: a
mismatched, unequal-mass grid (cells not of equal $\PP_0$-mass, though still
refining, with cell diameters $\to0$ and $\tau_S,\tau_Y$ a.e.\ continuous) is
expected, on the same mechanism as the proof above, to converge instead to the
correlation of $\tau_S,\tau_Y$ evaluated against the limiting distribution of cell
\emph{centers} rather than against $\PP_0$ itself, since Proposition~\ref{prop:interior}'s
reduction is unweighted over cells regardless of their true $\PP_0$-mass; we do not
prove this general claim here, but confirm it numerically for a specific
quantile-transform control grid, with $\tau_S(x)=x,\tau_Y(x)=x^5$ on
$\mathrm{Unif}[0,1]$: the dyadic grid targets $\sqrt{33}/7$ while the
quantile-transform grid targets $\sqrt3/2$, matching to within $3\times10^{-4}$ at
$K=256$. The equal-$\PP_0$-mass grid is the canonical choice
that makes the limit basis-independent and intrinsic to $\PP_0$; a different
(non-equal-mass) refining grid is not merely a slower route to the same answer, but
generally targets a different number entirely. \emph{Identification at each finite
level.} $\tau_a^{(m)}(k)=\E[\tau_a(X)\mid\text{cell }k]$
holds for the true CATEs by iterated expectations regardless of identification
strategy, but \emph{estimating} $\Theta_m$ from data at each level via the per-cell
AIPW/difference-of-means machinery of Section~\ref{sec:layer1} requires
unconfoundedness given the \emph{coarsened} covariate (the cell label) --- automatic
in a randomized trial, but not implied by unconfoundedness given the fine covariate
$X$ in an observational study, since coarsening a covariate can reintroduce
confounding. Proposition~\ref{prop:sieve-limit} is a statement about the estimand's
limit for the true CATEs; consistent estimation at each finite $m$ in an
observational study requires this additional, level-specific identification
condition.
\end{remark}

\subsection{Exiting the interior regime: a first-order account}\label{sec:exit-interior}

Propositions~\ref{prop:interior} and~\ref{prop:saturated} pin $\Theta(\PP_0,\lambda)$
down exactly at both edges of the $\lambda$-range; the intermediate range is where the
two closed forms need not hold (Remark~\ref{rem:saturated-scope}). This section
characterizes, in closed form, what happens \emph{just past} the interior edge
$\lambda=p_{\min}:=\min_kp_0(k)$: the sign and relative magnitude at which $\Theta$
first departs from $\Corr_K(\tau_S,\tau_Y)$ turn out to be governed by (a scalar
multiple of) the classical influence function of the Pearson correlation
coefficient, evaluated at the smallest covariate stratum's standardized CATEs --- an
MCMC-free determination of the \emph{sign} of that departure, computable directly
from $(\PP_0,\tau_S,\tau_Y)$, before any hit-and-run sampling.

\subsubsection{The interior body as a root polytope}

Write $H=\{V\in\R^K:\mathbf1^\top V=0\}$ for the zero-sum hyperplane,
$P_H:=I_K-\frac1K\mathbf1\mathbf1^\top$ for the orthogonal projection onto $H$ ---
the matrix already appearing, unnamed, as the factor of $a(\lambda)$ in
Proposition~\ref{prop:interior}'s proof (Appendix~\ref{app:proof-interior}) --- and
$\cB_\lambda:=\{V\in H:\|V\|_1\le2\lambda\}$ for the feasible set of
Section~\ref{sec:setup} before imposing nonnegativity. The actual feasible region is
$F_\lambda=\cB_\lambda\cap\{V:\PP_0+V\ge0\}$, equal to $\cB_\lambda$ exactly in the
interior regime ($\lambda\le p_{\min}$, Definition~\ref{def:interior}).

\begin{lemma}[Vertices and facets of $\cB_\lambda$]\label{lem:root-polytope}
For every $K\ge2$ and $\lambda>0$:
\begin{enumerate}[label=(\alph*)]
    \item $\cB_\lambda=\{V\in H:\sum_{k\in S}V_k\le\lambda\ \forall\,\emptyset\ne
    S\subsetneq[K]\}$ ($2^K-2$ facet inequalities, pairing $S$ with its complement
    $S^c$ to bound $|\sum_{k\in S}V_k|\le\lambda$ --- the standard
    $d_{\mathrm{TV}}=\sup_A|\Q(A)-\PP_0(A)|$ representation).
    \item $\cB_\lambda=\mathrm{conv}\{\lambda(e_i-e_j):1\le i\ne j\le K\}$, a polytope
    with exactly $K(K-1)$ vertices: the type-$A_{K-1}$ root polytope, scaled by
    $\lambda$ \citep[see][for the general theory of such
    bodies]{Postnikov2009}.
    \item For each cell $k$, the facet $\{V\in\cB_\lambda:V_k=-\lambda\}$ is itself
    the $(K-2)$-simplex $\lambda\cdot\mathrm{conv}\{e_i-e_k:i\ne k\}$; every other
    vertex has $k$-th coordinate $0$ or $+\lambda$.
\end{enumerate}
\end{lemma}

\begin{proof}[Proof sketch; full derivation in Appendix~\ref{app:proof-exit-interior}]
For $V\ne0$, (a) follows from the identity
\[
    \|V\|_1=2\max_{S\subseteq[K]}\sum_{k\in S}V_k, \qquad \text{valid whenever } \mathbf1^\top V=0
\]
(the maximizer is $S=\{k:V_k\ge0\}$, which by the zero-sum constraint carries
exactly half of $\|V\|_1$; $V=0$ is the midpoint of any antipodal vertex pair, so
needs no separate argument). (b) is a transportation argument: any $V\in\cB_\lambda$
decomposes as a nonnegative combination of the $K(K-1)$ two-cell moves
$\lambda(e_i-e_j)$, with coefficients summing to at most $1$ (feasibility of the
bipartite transportation problem matching $V$'s positive part to its negative
part), and each such point is extreme (an equality case of the triangle inequality
for $\|\cdot\|_1$). (c) follows because the $k$-th coordinate of a vertex
$\lambda(e_i-e_j)$ equals $-\lambda$ when $j=k$ and is otherwise $0$ (when
$i,j\ne k$) or $+\lambda$ (when $i=k$): $-\lambda$ is the unique minimum of the
linear functional $V\mapsto V_k$ over $\cB_\lambda$'s vertex set, so any convex
combination attaining it must place all its mass on the $K-1$ vertices
$\{\lambda(e_i-e_k)\}_{i\ne k}$ that themselves attain it.
\end{proof}

\begin{lemma}[Exact global constancy under a uniform $\PP_0$]\label{lem:uniform-global}
Suppose $\Var_{\PP_0}(\tau_S),\Var_{\PP_0}(\tau_Y)>0$
(Assumption~\ref{ass:nondegeneracy}'s finite-$K$ specialization). If
$p_0(k)=1/K$ for every $k=1,\dots,K$, then $\Theta(\PP_0,\lambda)=\Corr_K(\tau_S,\tau_Y)$
for \emph{every} $\lambda\in(0,1)$, not merely in the interior and saturated
regimes.
\end{lemma}

\begin{proof}
The nonnegativity region $\{V\in H:V_k\ge-1/K\ \forall k\}$ is invariant under every
permutation of $\{1,\dots,K\}$, for every $\lambda$: unlike the general case, this
does not depend on which cells are large or small, since every cell carries the same
mass. Hence $F_\lambda=\cB_\lambda\cap\{V:V_k\ge-1/K\ \forall k\}$ is
permutation-invariant at every $\lambda$, not only for $\lambda\le p_{\min}=1/K$, and
the permutation-symmetry argument of Proposition~\ref{prop:interior}'s proof ---
uniform measure on a permutation-invariant body forces $\Sigma=a(\lambda)P_H$, which
cancels exactly in the ratio $\Theta_{ab}=\tau_a^\top\Sigma\tau_b$ --- applies
verbatim for every $\lambda\in(0,1)$, not only for $\lambda\le p_{\min}$.
\end{proof}

\subsubsection{The first-order slope, and its influence-function identity}

\begin{proposition}[First-order exit from the interior regime]\label{prop:exit-interior}
Suppose $\Var_{\PP_0}(\tau_S),\Var_{\PP_0}(\tau_Y)>0$. Let $\cK^*=\arg\min_kp_0(k)$
and $p_{\min}=p_0(k)$ for $k\in\cK^*$. Write $r_K=\Corr_K(\tau_S,\tau_Y)$,
$c_a(k)=\tau_a(k)-\bar\tau_a$, and $z_a(k)=c_a(k)/\|c_a\|_2$ (the standardized,
centered CATE vectors already used in Proposition~\ref{prop:interior}'s proof). Then
$\Theta(\PP_0,\cdot)$ is right-differentiable at $\lambda=p_{\min}$, with
\begin{equation}\label{eq:exit-slope}
    \partial_\lambda^+\Theta(\PP_0,\lambda)\Big|_{\lambda=p_{\min}}
    = \frac{C_K}{p_{\min}}\sum_{k\in\cK^*}\varpi(k),
    \qquad
    \varpi(k) := \frac{r_K}2\big(z_S(k)^2+z_Y(k)^2\big) - z_S(k)z_Y(k),
\end{equation}
for a constant $C_K>0$ depending only on $K$
(Appendix~\ref{app:proof-exit-interior} expresses $C_K$ in terms of the root
polytope's facet-volume ratio, not evaluated in closed form here). Equivalently,
\begin{equation}\label{eq:exit-psi-alt}
    \varpi(k) = \frac{1+r_K}4\big(z_S(k)-z_Y(k)\big)^2 - \frac{1-r_K}4\big(z_S(k)+z_Y(k)\big)^2,
\end{equation}
proportional to the negative of the classical influence function of the Pearson
correlation coefficient \citep{DevlinGnanadesikanKettenring1975}, evaluated at cell
$k$'s standardized CATEs $(z_S(k),z_Y(k))$. Consequently, whenever
$\sum_{k\in\cK^*}\varpi(k)\ne0$, widening the ball past the interior regime raises
$\Theta$ above $\Corr_K$ if the smallest covariate strata are, in aggregate, a net
\emph{discordant} contributor to the baseline correlation, and lowers $\Theta$ if
they are, in aggregate, a net \emph{concordant} contributor; because $C_K$ is never
evaluated here, only the \emph{sign} of $\sum_{k\in\cK^*}\varpi(k)$ --- and ratios of
this quantity across designs or cells --- are available without further
calibration. Because a correlation's influence function integrates to zero over the
population, $\sum_{k=1}^K\varpi(k)=0$; when every cell ties at the minimum
($p_0(1)=\dots=p_0(K)=1/K$), \eqref{eq:exit-slope}'s right-hand side vanishes
identically, consistent with (though not, by itself, a proof of)
Lemma~\ref{lem:uniform-global}'s stronger, independently-established claim that
$\Theta\equiv\Corr_K$ throughout $(0,1)$ in that case.
\end{proposition}

\begin{remark}[Scope: a local result, not a global trajectory]\label{rem:exit-scope}
Proposition~\ref{prop:exit-interior} characterizes the one-sided derivative exactly
at $\lambda=p_{\min}$; by itself it says nothing about $\Theta$'s trajectory deeper
in the intermediate regime $(p_{\min},1-p_{\min})$, where later combinatorial events
(Remark~\ref{rem:kink-subset-sums}), not captured by this first-order term,
generally take over. The sign prediction requires no calibration of $C_K$ and is
validated numerically across four designs (one with a unique minimizer and
comfortably large $p_{\min}=0.10$; two with a unique minimizer and
$p_{\min}\approx0.04$--$0.05$ constructed to produce opposite signs; and the
paper's own canonical DGP~1, whose minimizer is tied by symmetry and whose
predicted effect is the smallest of the four): a hit-and-run simulation with common
random numbers across $\lambda$, resolving $\lambda\in(p_{\min},1.3p_{\min})$
specifically to stay inside the window where the first-order approximation should
dominate, confirms the predicted sign in every case, including DGP~1's small,
tied-minimizer effect. Fitting a single positive constant $\hat C_K$ to the four
measured initial slopes against $\sum_{k\in\cK^*}\varpi(k)/p_{\min}$ (itself a
further validation, since a negative or inconsistent fit would refute the
proposition rather than merely fail to calibrate it) gives $\hat C_K>0$ with the
four implied values of $C_K$ agreeing to within $30\%$ of each other despite the
designs' very different $p_{\min}$, tie structure, and CATE patterns. Deeper in the
band --- e.g.\ near the peak $\lambda\approx0.25$ observed for DGP~1 in the
validation study (Example~\ref{ex:binding}), several subset-sum thresholds past
$p_{\min}$, where the sign of the trajectory need not match this local prediction
--- this proposition does not apply, and the trajectory is genuinely empirical, not
a corollary of this local result.
\end{remark}

\begin{remark}[The kink structure is governed by subset sums, not individual cell
masses]\label{rem:kink-subset-sums}
Cell $k$'s box constraint $V_k\ge-p_0(k)$ first binds at $\lambda=p_0(k)$, but once
several cells' constraints are simultaneously active, a \emph{pair} $\{i,j\}$ can be
jointly driven toward zero only once $\lambda\ge p_0(i)+p_0(j)$, and analogously for
larger subsets: the combinatorial type of $F_\lambda$ changes exactly at the values
$\lambda\in\{p_0(S):\emptyset\ne S\subsetneq[K]\}$, the \emph{subset sums} of
$\PP_0$'s cell probabilities (up to $2^K-2$ distinct values, not $K$). The two band
endpoints of Propositions~\ref{prop:interior} and~\ref{prop:saturated},
$p_{\min}=\min_kp_0(k)$ and $1-p_{\min}=\max_{S\subsetneq[K]}p_0(S)$, are exactly the
two extreme subset sums; Proposition~\ref{prop:exit-interior} characterizes only the
first such event, at $\lambda=p_{\min}$.
\end{remark}


% ===========================================================================
\section{An alternative geometry: $\chi^2$ balls}\label{sec:chisq}
% ===========================================================================

Sections~\ref{sec:interior}--\ref{sec:exit-interior} develop the theory for
$d=d_{\mathrm{TV}}$, the paper's primary instantiation of the generic local
geometry $U(\PP_0,\lambda;d)$ of Section~\ref{sec:setup}. This section shows that
TV is not \emph{uniquely required} for an exact, $\lambda$-free interior-regime
closed form to exist: $\chi^2$ divergence,
\begin{equation}\label{eq:chisq-def}
    \chi^2(\Q,\PP_0) := \sum_{k=1}^Kp_0(k)\Big(\frac{q_k}{p_0(k)}-1\Big)^2
    = \sum_{k=1}^K\frac{(q_k-p_0(k))^2}{p_0(k)} = \big\|\tfrac{d\Q}{d\PP_0}-1\big\|_{L^2(\PP_0)}^2,
\end{equation}
admits its own such closed form --- but the closed form it yields is specific to
$\chi^2$, not the one TV yields: the two geometries' interior-regime values
\emph{differ} ($\Corr_{\PP_0}$ versus $\Corr_K$), even though their saturated-regime
values coincide, via a mechanism (Section~\ref{sec:saturated}) that turns out to be
distance-free. Throughout this section, $\PP_0$ has full support ($p_0(k)>0$ for
every $k=1,\dots,K$ --- needed for $D:=\mathrm{diag}(p_0)$ to be invertible, and
already the standing convention of Section~\ref{sec:setup}, where cells with
$p_0(k)=0$ are excluded from the geometry entirely) and
$\Var_{\PP_0}(\tau_S),\Var_{\PP_0}(\tau_Y)>0$ (Assumption~\ref{ass:nondegeneracy})
as before. $\mu_\lambda$ continues to denote the measure uniform with respect to
Lebesgue measure on the zero-sum hyperplane $H$, restricted to $U(\PP_0,\lambda;d)$:
this reference-measure convention, held fixed across $d$, is what makes the
comparison below a theorem about $(d,\mu)$ pairs rather than an artifact of also
changing what ``uniform'' means.

\subsection{The $\chi^2$-ball interior closed form}

\begin{definition}[$\chi^2$ interior regime]\label{def:interior-chisq}
The $\chi^2$-ball $U(\PP_0,\lambda;\chi^2)$ is in the \emph{interior regime} if the
nonnegativity constraint does not bind: $\{V\in H:V^\top D^{-1}V\le\lambda\}\subseteq
\{V:\PP_0+V\ge0\}$, paralleling Definition~\ref{def:interior} for TV.
\end{definition}

\begin{proposition}[$\chi^2$ interior-regime closed form]\label{prop:interior-chisq}
For any $\lambda>0$ in the interior regime of Definition~\ref{def:interior-chisq},
\begin{equation}\label{eq:interior-chisq}
    \Theta(\PP_0,\lambda) = \Corr_{\PP_0}(\tau_S,\tau_Y)
    := \frac{\Cov_{\PP_0}(\tau_S,\tau_Y)}{\sqrt{\Var_{\PP_0}(\tau_S)\Var_{\PP_0}(\tau_Y)}},
\end{equation}
the ordinary $\PP_0$-weighted correlation of the CATEs --- the same object
Proposition~\ref{prop:sieve-limit} identifies as a continuous-$X$ \emph{limit} on an
equal-$\PP_0$-mass dyadic grid, but here exact on \emph{any} finite $\cX$, at
\emph{every} $\lambda$ in the interior regime, and with no equal-mass grid
requirement: it removes, for $\chi^2$, the grid-dependence caveat
Remark~\ref{rem:sieve-scope} documents for TV's continuous-$X$ limit. In particular
$\Theta(\PP_0,\lambda)$ does not depend on $\lambda$.
\end{proposition}

\begin{proof}
Substitute $G=D^{-1/2}V$, i.e.\ $V=D^{1/2}G$ ($D$ is positive definite by the
standing full-support assumption). Since this is a linear bijection between the
$(K-1)$-dimensional subspaces $H$ and $H':=\{G:\langle\nu,G\rangle=0\}$,
$\nu:=D^{1/2}\mathbf1=(\sqrt{p_0(k)})_k$ (using
$\mathbf1^\top V=\mathbf1^\top D^{1/2}G=\langle\nu,G\rangle$), it scales
$(K-1)$-dimensional Hausdorff measure by a constant Jacobian factor (computed in
orthonormal bases of $H$ and $H'$), hence carries the uniform law on any body in
$H$ to the uniform law on its image in $H'$; $\nu$ is a \emph{unit} vector, since
$\|\nu\|^2=\sum_kp_0(k)=1$. The quadratic constraint becomes
$V^\top D^{-1}V=G^\top D^{1/2}D^{-1}D^{1/2}G=\|G\|_2^2\le\lambda$: in the interior
regime the feasible body for $G$ is exactly the Euclidean ball
$\{G\in H':\|G\|_2\le\sqrt\lambda\}$, invariant under the stabilizer of $\nu$ in
$O(K)$, which acts on $H'$ as the full orthogonal group $O(K-1)$. By this
invariance, $\Cov(G)$ restricted to $H'$ is a scalar multiple of the identity there,
$\Cov(G)=c(I_K-\nu\nu^\top)$; the standard fact that uniform measure on a
$d$-dimensional ball of radius $r$ has covariance $\frac{r^2}{d+2}I_d$ (here
$d=K-1$, $r=\sqrt\lambda$) gives $c=\lambda/(K+1)$, so
$\Cov(G)=\frac\lambda{K+1}(I_K-\nu\nu^\top)$ as a $K\times K$ matrix. Un-whitening with $D^{1/2}\nu=D\mathbf1=p_0$ (the vector of cell
probabilities),
\[
    \Sigma=\Cov(V)=D^{1/2}\Cov(G)D^{1/2}
    =\frac\lambda{K+1}\big(D-(D^{1/2}\nu)(D^{1/2}\nu)^\top\big)
    =\frac\lambda{K+1}\big(D-p_0p_0^\top\big).
\]
Substituting into $\Theta_{ab}=\tau_a^\top\Sigma\tau_b$,
$\Theta_{ab}=\frac\lambda{K+1}\big(\tau_a^\top D\tau_b-(\tau_a^\top p_0)(\tau_b^\top
p_0)\big)=\frac\lambda{K+1}\Cov_{\PP_0}(\tau_a,\tau_b)$; forming the ratio
$\Theta_{SY}/\sqrt{\Theta_{SS}\Theta_{YY}}$ cancels $\lambda/(K+1)$ exactly, giving
\eqref{eq:interior-chisq} --- the same degree-zero-homogeneity cancellation as
Proposition~\ref{prop:interior}'s $a(\lambda)$.
\end{proof}

\begin{lemma}[Sharp $\chi^2$ interior threshold]\label{lem:interior-chisq-threshold}
The $\chi^2$-ball's interior regime (Definition~\ref{def:interior-chisq}) holds for
$\lambda\le p_{\min}/(1-p_{\min})$, $p_{\min}:=\min_kp_0(k)$, and this threshold is
sharp. On a continuous or countably-infinite-atom $\cX$ (even with full support in
the sense that every open neighborhood has positive $\PP_0$-mass), the interior
regime is empty at every $\lambda>0$, by the same mechanism as
Lemma~\ref{lem:interior-atomic}: bounding $\chi^2(\Q,\PP_0)$ controls
$\|d\Q/d\PP_0-1\|_{L^2(\PP_0)}$, not $\|d\Q/d\PP_0-1\|_{L^\infty(\PP_0)}$, so a
small-mass event can still be driven arbitrarily close to zero probability at
arbitrarily small $\chi^2$ cost.
\end{lemma}

\begin{proof}[Proof sketch]
For fixed $k$, minimizing $V^\top D^{-1}V=\sum_jV_j^2/p_0(j)$ subject to
$V_k=-p_0(k)$ and $\sum_jV_j=0$ (i.e.\ $\sum_{j\ne k}V_j=p_0(k)$) is, by Lagrange
multipliers, solved by $V_j=p_0(j)p_0(k)/(1-p_0(k))$ for $j\ne k$ (proportional
redistribution minimizes a $D$-weighted sum of squares under a linear constraint),
giving minimum cost $\sum_{j\ne k}V_j^2/p_0(j)+V_k^2/p_0(k)=p_0(k)^2/(1-p_0(k))+p_0(k)
=p_0(k)/(1-p_0(k))$: cell $k$'s odds. Since $x\mapsto x/(1-x)$ is increasing on
$[0,1)$, the overall threshold --- the smallest cost at which \emph{any} cell's
constraint first binds --- is $\min_kp_0(k)/(1-p_0(k))=p_{\min}/(1-p_{\min})$, and
this is exactly attained (not merely a bound), by the same argument
Lemma~\ref{lem:interior-atomic} uses for TV. For the continuous-$X$ failure, mirror
Lemma~\ref{lem:interior-atomic}'s ($\Rightarrow$) direction exactly: for any
$\lambda>0$, choose $A$ with $0<\PP_0(A)=\varepsilon$ small enough that
$\varepsilon/(1-\varepsilon)<\lambda$ (possible since the left side $\to0$ as
$\varepsilon\to0$). The perturbation $v=-\mathbf1_A+\frac\varepsilon{1-\varepsilon}\mathbf1_{A^c}$
drives the density to exactly $0$ on $A$ at $\chi^2$ cost
$\varepsilon^2/\varepsilon+\big(\frac\varepsilon{1-\varepsilon}\big)^2(1-\varepsilon)=\varepsilon+\varepsilon^2/(1-\varepsilon)
=\varepsilon/(1-\varepsilon)$, strictly inside the budget $\lambda$; rescaling $v$ by a factor
$1+\delta$ for small enough $\delta>0$ (remaining feasible, since the $\chi^2$ cost
scales as $(1+\delta)^2\varepsilon/(1-\varepsilon)$, still $<\lambda$ for
$\delta<\sqrt{\lambda(1-\varepsilon)/\varepsilon}-1>0$) drives the density \emph{strictly
negative} on $A$, violating nonnegativity within the $\chi^2$ budget $\lambda$. So
no $\lambda>0$ excludes every such $A$ once $\PP_0$ has a non-atomic component or
infinitely many atoms.
\end{proof}


\subsection{The $\chi^2$ instance of saturation}

Lemma~\ref{lem:saturation-general} (Section~\ref{sec:saturated}) already settles
saturation for \emph{any} distance, including $\chi^2$; what remains is to compute
the $\chi^2$-specific radius at which $U(\PP_0,\lambda;\chi^2)$ reaches the whole
simplex.

\begin{corollary}[$\chi^2$ instance of Lemma~\ref{lem:saturation-general}]\label{cor:saturation-chisq}
Suppose $p_0(k)>0$ for all $k$ (full support). Then $\chi^2(\cdot,\PP_0)$ is a
positive-definite quadratic form in $\Q-\PP_0$ (from \eqref{eq:chisq-def},
$D^{-1}\succ0$), hence strictly convex in $\Q$, so its maximum over $\Delta_{K-1}$
is attained at a vertex $\delta_{k^*}$:
$\chi^2(\delta_{k^*},\PP_0)=\frac{(1-p_0(k^*))^2}{p_0(k^*)}+(1-p_0(k^*))
=\frac{1-p_0(k^*)}{p_0(k^*)}$, maximized (since $(1-x)/x$ is decreasing) at
$k^*=\arg\min_kp_0(k)$: $U(\PP_0,\lambda;\chi^2)=\Delta_{K-1}$ for
$\lambda\ge(1-p_{\min})/p_{\min}$, the exact reciprocal of
Lemma~\ref{lem:interior-chisq-threshold}'s interior threshold. By
Lemma~\ref{lem:saturation-general}, $\Theta(\PP_0,\lambda)=\Corr_K(\tau_S,\tau_Y)$
there --- the same value Corollary~\ref{cor:saturation-tv} below gives for TV, even
though the two geometries' \emph{interior} values differ
(Remark~\ref{rem:tv-chisq-comparison}). If $\PP_0$ has any cell with $p_0(k)=0$,
$\chi^2(\delta_k,\PP_0)=\infty$ for that $k$ and the $\chi^2$ ball never saturates ---
unlike TV, which always saturates by Lemma~\ref{lem:saturation} regardless of
$\PP_0$'s support. (The reverse convention $\sum_k(p_0(k)-q_k)^2/q_k$ is likewise
$+\infty$ at every vertex and never saturates; \eqref{eq:chisq-def}'s convention,
with $\PP_0$ in the denominator, is the one used throughout.)
\end{corollary}

\begin{corollary}[TV instance, restated]\label{cor:saturation-tv}
Lemma~\ref{lem:saturation} and Proposition~\ref{prop:saturated} are the
$d=d_{\mathrm{TV}}$ instance of Lemma~\ref{lem:saturation-general}: saturation at
$\lambda\ge1-p_{\min}$, with $\Theta(\PP_0,\lambda)=\Corr_K(\tau_S,\tau_Y)$ there.
\end{corollary}

\begin{remark}[Two geometries, compared]\label{rem:tv-chisq-comparison}
\begin{center}
\begin{tabular}{@{}lcccc@{}}
\toprule
& interior threshold $\lambda^*$ & $\Theta$ at $\lambda\le\lambda^*$
  & saturation $\lambda_{\mathrm{sat}}$ & $\Theta$ at $\lambda\ge\lambda_{\mathrm{sat}}$ \\
\midrule
TV & $p_{\min}$ & $\Corr_K(\tau_S,\tau_Y)$ & $1-p_{\min}$ & $\Corr_K(\tau_S,\tau_Y)$ \\
$\chi^2$ & $p_{\min}/(1-p_{\min})$ & $\Corr_{\PP_0}(\tau_S,\tau_Y)$
  & $(1-p_{\min})/p_{\min}$ & $\Corr_K(\tau_S,\tau_Y)$ \\
\bottomrule
\end{tabular}
\end{center}
(The $\lambda^*,\lambda_{\mathrm{sat}}$ columns are not on a common numeric scale
across rows; see the Cauchy--Schwarz calibration below before comparing them.)
Under TV, $\lambda^*+\lambda_{\mathrm{sat}}=1$ (complementary \emph{probabilities});
under $\chi^2$, $\lambda^*\cdot\lambda_{\mathrm{sat}}=1$ (reciprocal \emph{odds}) ---
two computations specific to each geometry, not instances of one general identity.
Both endpoints coincide for TV, so the TV radius carries no information about which
of two named estimands is being approached; under $\chi^2$, by contrast, $\lambda$
genuinely interpolates between two distinct, recognizable quantities
($\Corr_{\PP_0}\to\Corr_K$) --- a $\chi^2$-specific advantage for interpreting what
widening the ball means, at the cost of $\chi^2$ never saturating unless $\PP_0$ has
full support. The two $\lambda$-scales are not on the same numeric footing: by
Cauchy--Schwarz,
\[
    2d_{\mathrm{TV}}(\Q,\PP_0)=\sum_k\sqrt{p_0(k)}\cdot\frac{|q_k-p_0(k)|}{\sqrt{p_0(k)}}
    \le\sqrt{\textstyle\sum_kp_0(k)}\cdot\sqrt{\chi^2(\Q,\PP_0)}=\sqrt{\chi^2(\Q,\PP_0)},
\]
so a $\chi^2$-ball of radius $\lambda_{\chi^2}$ is \emph{contained} in the TV ball of
radius $\tfrac12\sqrt{\lambda_{\chi^2}}$: this is a one-directional containment,
not a tight numerical equivalence (e.g.\ at $\lambda_{\chi^2}=\lambda^*_{\chi^2}=
p_{\min}/(1-p_{\min})=0.0526$ for DGP1's $p_{\min}=0.05$, the bound gives a TV
radius of $0.115$, larger than TV's own interior threshold $p_{\min}=0.05$ --- a
loose bound, not a contradiction), and the table's $\lambda_{\chi^2}$ column should
not be read on the same scale as $\lambda_{\mathrm{TV}}$ without it. Finally: the
intermediate-regime numbers already reported in Example~\ref{ex:binding} --- the
unweighted $\Corr_K=0.6455$ and the ``$p_0$-weighted $0.6411$'' --- are two
\emph{different geometries'} interior (and, for the first, also saturated) closed
forms, both evaluated at DGP1's $\tau_S,\tau_Y,\PP_0$: $0.6455$ is the TV
geometry's closed form \eqref{eq:interior-reduction}, and $0.6411$ is exactly the
$\chi^2$-ball's interior closed form $\Corr_{\PP_0}(\tau_S,\tau_Y)$ of
\eqref{eq:interior-chisq}, which this section identifies as such for the first
time. Neither is $\Theta_{\chi^2}$ or $\Theta_{\mathrm{TV}}$ evaluated at DGP1's
validated $\lambda=0.3$, which lies well outside \emph{both} geometries' interior
regimes ($\lambda^*_{\chi^2}=0.05/0.95\approx0.053$, $\lambda^*_{\mathrm{TV}}=0.05$);
\eqref{eq:interior-chisq} identifies what the $0.6411$ \emph{is}, not a new
computation at $\lambda=0.3$.
\end{remark}

\begin{remark}[Scope: an exact result for two geometries, an asymptotic one for the
rest]\label{rem:chisq-general}
The mechanism generating Proposition~\ref{prop:interior-chisq} is not unique to
$\chi^2$ among divergences from $\PP_0$: for any $f$-divergence
$D_f(\Q\|\PP_0)=\sum_kp_0(k)f(q_k/p_0(k))$ with $f(1)=0$ and $f$ twice
differentiable at $1$ with $f''(1)\in(0,\infty)$, the linear term in $V$ vanishes on
$H$ (since $\sum_kV_k=0$), giving $D_f=\tfrac12f''(1)\chi^2(\Q,\PP_0)+o(\|V\|^2)$,
the remainder uniform in $k$ (hence summable against the finitely many, boundedly
away from $0$, cell probabilities $p_0(k)$) since $K$ is finite and $p_0$ has full
support: to second order near $\PP_0$, every such $D_f$-ball is the same $\chi^2$
ellipsoid up to an overall scale, which cancels in $\Theta$'s ratio exactly as
$a(\lambda)$ and $\lambda/(K+1)$ do above. This suggests
$\Theta_f(\lambda)\to\Corr_{\PP_0}(\tau_S,\tau_Y)$
as $\lambda\downarrow0$ for KL divergence, squared Hellinger distance, and other
common choices sharing this local quadratic structure --- with total variation the
one common exception, since $f(t)=\tfrac12|t-1|$ has a kink at $t=1$ rather than a
well-defined second derivative there, which is exactly why its interior geometry is
the polytope of Lemma~\ref{lem:root-polytope} rather than an ellipsoid. This is a
weaker, asymptotic ($\lambda\downarrow0$ only) claim than
Proposition~\ref{prop:interior-chisq}'s exact, $\lambda$-free statement for
$\chi^2$ itself --- $D_f$ does not equal $\chi^2$ exactly for $\lambda$ bounded away
from $0$, so its interior regime need not be $\lambda$-free the way $\chi^2$'s and
TV's are --- and is left as a direction for future work rather than developed here.
\end{remark}


% ===========================================================================
\section{Finite-support theory}\label{sec:layer1}
% ===========================================================================

Fix a finite (or discretized) covariate space $\cX=\{1,\dots,K\}$ with cell
probabilities $p_0(k)=\PP_0(X=k)$. We observe treatment $A$, surrogate $S$, and
outcome $Y$; write $\mu_a^S(k)=\E[S\mid A=a,X=k]$, $\mu_a^Y(k)=\E[Y\mid A=a,X=k]$, and
recall $\tau_S(k)=\mu_1^S(k)-\mu_0^S(k)$, $\tau_Y(k)=\mu_1^Y(k)-\mu_0^Y(k)$ from
\eqref{eq:cate}. For a future study with cell masses $q=(q_1,\dots,q_K)$,
\eqref{eq:delta-linear} specializes to
\begin{equation}\label{eq:delta-discrete}
    \Delta_S(\Q) = \sum_k q_k\,\tau_S(k), \qquad \Delta_Y(\Q) = \sum_k q_k\,\tau_Y(k),
\end{equation}
and \eqref{eq:discrete-quadform} gives the estimand as the quadratic form
$\Theta_{SY}=\tau_S^\top\Sigma\tau_Y$ normalized by $\sqrt{\Theta_{SS}\Theta_{YY}}$.

\subsection{Estimator}

Given $\Q_1,\dots,\Q_M$ sampled by hit-and-run from $U(\Pn,\lambda;d_{\mathrm{TV}})$
(Convention~\ref{conv:geometry}), and importance weights $w_i(\Q_m)=q_{m,k}/\hat
p_0(k)$ for $X_i=k$, the cross-fitted AIPW estimator of $\Delta_S(\Q_m)$ is
\begin{equation}\label{eq:aipw-estimator}
    \hat\Delta_S(\Q_m) = \frac1n\sum_{i=1}^n w_i(\Q_m)\left[
        \frac{A_i(S_i-\hat\mu_1^S(X_i))}{\ehat(X_i)}
        - \frac{(1-A_i)(S_i-\hat\mu_0^S(X_i))}{1-\ehat(X_i)}
        + \hat\mu_1^S(X_i)-\hat\mu_0^S(X_i)
    \right],
\end{equation}
with $\hat\Delta_Y(\Q_m)$ defined analogously, and nuisances $\ehat,\hat\mu_a$
cross-fitted (Assumption~\ref{ass:nuisance}). The plug-in estimator of the estimand is
the sample Pearson correlation of the $M$ pairs,
\begin{equation}\label{eq:theta-hat}
    \Thetahat_n(\lambda) = \Corr\!\big(\{\hat\Delta_S(\Q_m)\}_{m=1}^M,\{\hat\Delta_Y(\Q_m)\}_{m=1}^M\big).
\end{equation}

\begin{lemma}[Exact cancellation of the ball's empirical center]\label{lem:cancellation}
Write $\hat p_0(k)=n_k/n$ for the empirical cell frequency, $n_k=\sum_i\mathbf1\{X_i=k\}$.
Then, for any fixed weight vector $q_m$,
\begin{equation}\label{eq:cancellation}
    \hat\Delta_S(\Q_m) = \sum_{k=1}^K q_{m,k}\,\hat\tau_S(k),
    \qquad
    \hat\tau_S(k) = \frac1{n_k}\sum_{i:X_i=k}\left[
        \frac{A_i(S_i-\hat\mu_1^S(X_i))}{\ehat(X_i)}
        - \frac{(1-A_i)(S_i-\hat\mu_0^S(X_i))}{1-\ehat(X_i)}
        + \hat\mu_1^S(X_i)-\hat\mu_0^S(X_i)
    \right],
\end{equation}
an identity that holds for every finite sample, independently of $\hat p_0$. In
particular, using $\Pn$ rather than $\PP_0$ to define the importance weights
$w_i(\Q_m)=q_{m,k}/\hat p_0(k)$ introduces no error, of any order, into
$\hat\Delta_S(\Q_m)$ conditional on $\Q_m$; $\hat\tau_S(k)$ is exactly the per-cell
cross-fitted AIPW estimator of $\tau_S(k)$.
\end{lemma}

\subsection{Main theorem}

\begin{theorem}[Correlation of treatment effects across studies]\label{thm:main}
Under Assumptions~\ref{ass:iid}--\ref{ass:geometry}, with $\Thetahat_n(\lambda)$ as in
\eqref{eq:theta-hat}, there is a mean-zero influence function
$\psi_\Theta(\cdot;\Qhat)$ with $\E_{\PP_0}[\psi_\Theta(O;\Qhat)^2]=\sigma^2(\Qhat)<\infty$
such that:
\begin{enumerate}[label=(\alph*)]
    \item \textbf{Conditional statement (any $M$).} Conditional on
    $\{\Q_1,\dots,\Q_M\}$, as $n\to\infty$,
    \begin{equation}\label{eq:thm-main-conditional}
        \sqrt n\{\Thetahat_n(\lambda) - \Theta(\Qhat,\lambda)\}
        \xrightarrow{d} N\big(0,\sigma^2(\Qhat)\big),
    \end{equation}
    where $\Theta(\Qhat,\lambda)$ is the correlation functional evaluated at the
    empirical measure $\Qhat=\frac1M\sum_m\delta_{\Q_m}$ of the sampled studies.
    \item \textbf{Unconditional statement (if $n=o(M)$).} If additionally $M=M(n)\to\infty$
    with $n/M\to0$,
    \begin{equation}\label{eq:thm-main-unconditional}
        \sqrt n\{\Thetahat_n(\lambda) - \Theta(\PP_0,\lambda)\}
        \xrightarrow{d} N\big(0,\sigma^2(\mu)\big).
    \end{equation}
    (The direction is $n=o(M)$: $M$ must grow \emph{faster} than $n$, not slower.)
    \item \textbf{Combined statement (any $n,M$).}
    \begin{equation}\label{eq:thm-main-combined}
        \Thetahat_n(\lambda) - \Theta(\PP_0,\lambda) = O_P(n^{-1/2}) + O_P(M^{-1/2}),
    \end{equation}
    with the two error sources asymptotically independent (the first is a function of
    $O_1,\dots,O_n$ alone, the second of the Markov chain alone, conditional on which
    the first is computed); a standard error accounting for both is
    $\widehat{\mathrm{se}} = \sqrt{\hat\sigma^2(\Qhat)/n + \hat\kappa^2/M}$, with
    $\hat\kappa^2$ a batch-means or spectral estimate of the Monte Carlo asymptotic
    variance of Lemma~\ref{lem:mcmc-error}.
\end{enumerate}
The influence function is
\begin{equation}\label{eq:psi-theta-layer1}
    \psi_\Theta(O;\Qhat) = \frac1M\sum_{m=1}^M \nabla\varphi(\hat m)^\top
    \begin{pmatrix}\psi_S^{\mathrm{AIPW}}(O;\Q_m)\\ \psi_Y^{\mathrm{AIPW}}(O;\Q_m)\end{pmatrix},
\end{equation}
where $\varphi$ is the correlation map applied to the vector of five across-study
moments $m=(\E_\mu\Delta_S,\E_\mu\Delta_Y,\E_\mu\Delta_S^2,\E_\mu\Delta_Y^2,\E_\mu\Delta_S\Delta_Y)$,
and, for each fixed $\Q_m$ with weights $w(\cdot;\Q_m)$,
\begin{equation}\label{eq:psi-aipw}
    \psi_S^{\mathrm{AIPW}}(O;\Q_m) = w(X;\Q_m)\left[
        \frac{A(S-\mu_1^S(X))}{e(X)} - \frac{(1-A)(S-\mu_0^S(X))}{1-e(X)}
        + \mu_1^S(X)-\mu_0^S(X)
    \right] - \Delta_S(\Q_m),
\end{equation}
with $\psi_Y^{\mathrm{AIPW}}$ defined analogously. \emph{The centering is the constant
$-\Delta_S(\Q_m)$, subtracted outside the importance weight}, not
$-w(X;\Q_m)\Delta_S(\Q_m)$: only constant centering gives
$\E_{\PP_0}[\psi_S^{\mathrm{AIPW}}(O;\Q_m)]=0$ (Lemma~\ref{lem:if-meanzero},
Appendix~\ref{app:proofs}). The whole bracketed AIPW score, including the plug-in
regression term, is multiplied by $w(X;\Q_m)$: this differs from the textbook
doubly-robust average-treatment-effect influence function (which has an additional,
un-reweighted covariate-law augmentation term) because the target
$\Delta_S(\Q_m)=\E_{\Q_m}[\mu_1^S(X)-\mu_0^S(X)]$ is an expectation under the
\emph{fixed, external} composition $\Q_m$, not under $\PP_0$, so there is no
$\PP_0$-covariate-law term to augment.
\end{theorem}

\begin{remark}[Practical use]\label{rem:practical-use}
In the four-DGP validation study, $M\approx2100\gg n^{1/2}$ but $n=10{,}000\not=o(M)$
in the sense of part (b): the study does not drive $M/n\to\infty$, so the
\emph{conditional} statement (a) is the one the software's reported standard error
targets ($\Theta(\Qhat,\lambda)$, not $\Theta(\PP_0,\lambda)$), and the appropriate
label for reported coverage is coverage of $\Theta(\Qhat,\lambda)$. The combined
statement (c) is what practice should ultimately report; under (a)'s stated regime,
part (c)'s two terms are of comparable order only when $M\asymp n$, and the
Monte-Carlo term is, in the validation study's regime, small relative to sampling
variability but not derived to be asymptotically negligible without invoking (b).
\end{remark}

% ===========================================================================
\section{The randomized-trial corollary}\label{sec:rct}
% ===========================================================================

Theorem~\ref{thm:main} is stated for the cross-fitted AIPW estimator. The
randomized-trial, importance-weighting estimator used elsewhere in the package
(\texttt{method = "importance\_weighting"}) is not a separately-derived special case;
it is the AIPW estimator under two simultaneous degeneracies, and its influence
function follows by exact algebraic substitution.

\begin{corollary}[Randomized-trial reduction]\label{cor:rct}
Suppose (i) $A\perp X$ under $\PP_0$ with known constant propensity $e(X)\equiv\pi$
(randomization), and (ii) the working outcome regressions are taken to be constants
within arm, $\mu_a^S(X)\equiv m_{S,a}$ (no covariate adjustment). Then $\ehat=\bar
e_1:=\frac1n\sum_iw_i(\Q_m)A_i$ (the weighted arm-1 proportion), and $\bar e_1+\bar
e_0=1$ automatically since $\E_{\PP_0}[w]=1$ (there is no separate estimate of
$\bar e_0$, and no factor of $2$ anywhere unless $w\equiv1,\pi=1/2$). Taking
$\hat\mu_a^S=\hat m_{S,a}=\sum_iw_i(\Q_m)A_i S_i/\sum_iw_i(\Q_m)A_i$ (so that
$\sum_iw_i(\Q_m)A_i(S_i-\hat m_{S,1})=0$ identically), the AIPW estimator
\eqref{eq:aipw-estimator} collapses \emph{exactly} to the Hájek weighted
difference-in-means,
\begin{equation}\label{eq:iw-estimator}
    \hat\Delta_S(\Q_m) = \hat m_{S,1}-\hat m_{S,0}
    = \frac{\sum_i w_i(\Q_m)A_iS_i}{\sum_i w_i(\Q_m)A_i}
    - \frac{\sum_i w_i(\Q_m)(1-A_i)S_i}{\sum_i w_i(\Q_m)(1-A_i)},
\end{equation}
with influence function
\begin{equation}\label{eq:psi-iw}
    \psi_S^{\mathrm{IW}}(O;\Q_m) = w(X;\Q_m)\left[
        \frac{A(S-m_{S,1})}{\bar e_1} - \frac{(1-A)(S-m_{S,0})}{\bar e_0}
    \right],
\end{equation}
where $\bar e_1=\E[wA]$, $\bar e_0=\E[w(1-A)]$ are the average weights per arm
(estimated by $\frac1n\sum_iw_iA_i$, $\frac1n\sum_iw_i(1-A_i)$). Theorem~\ref{thm:main}
applies verbatim with $\psi_S^{\mathrm{IW}}$ in place of $\psi_S^{\mathrm{AIPW}}$.
\end{corollary}

\begin{remark}[Why this is a genuine special case, not a re-derivation]\label{rem:rct-scope}
Both degeneracies matter: randomization alone (known constant $e$) does \emph{not}
collapse AIPW to the weighted difference in means when $\mu_a^S(X)$ is a genuine
covariate-specific model, because covariate-adjusted estimators are strictly more
efficient than unadjusted ones even under randomization; only forcing the working
outcome model to be arm-constant (no adjustment) makes the IPW-residual terms vanish
identically. This is exactly what the package's \texttt{"importance\_weighting"}
method implements (no outcome regression step at all), and exactly why
Theorem~\ref{thm:main}'s AIPW estimator, with a real covariate-adjusted outcome model,
is a distinct (and more efficient) estimator even in a randomized trial.
\end{remark}

% ===========================================================================
\section{General theory: a bilinear-functional estimand}\label{sec:layer2}
% ===========================================================================

We now drop the finite-support restriction, allowing $X$ to be continuous on
$\cX\subseteq\R^d$, and generalize the target from the correlation
\eqref{eq:estimand} to any smooth function of finitely many across-study moments.
This section is not implemented in the package and not exercised by the simulation
study (Section~\ref{sec:scope}); it is included because it strictly generalizes
Section~\ref{sec:layer1}, verified by the recovery corollary of
Section~\ref{sec:recovery}.

\subsection{The estimand as a smooth function of linear and bilinear CATE functionals}

Let $\Psi:\R^K\to\R$ be smooth, and let
\begin{equation}\label{eq:general-estimand}
    \Theta = \Psi(m_1,\dots,m_K), \qquad m_j = \E_{\Q\sim\mu}[\phi_j(\Delta_S(\Q),\Delta_Y(\Q))],
\end{equation}
for base moments $\phi_j$. Correlation, $R^2$, mean squared prediction error, and
concordance are instances, all expressible through the five moments
$m=(\E_\mu\Delta_S,\E_\mu\Delta_Y,\E_\mu\Delta_S^2,\E_\mu\Delta_Y^2,\E_\mu\Delta_S\Delta_Y)$.
By \eqref{eq:reduction-linear}--\eqref{eq:reduction-bilinear}, each $m_j$ is built from
the linear functionals $\E_\mu\Delta_a=\int\bar a\,\tau_a\,d\PP_0$ and the bilinear
functionals $\Theta_{ab}=\iint C(x,x')\tau_a(x)\tau_b(x')\,d\PP_0(x)\,d\PP_0(x')$,
$a,b\in\{S,Y\}$ (using $\E_\mu[\Delta_a^2]=\Theta_{aa}+(\E_\mu\Delta_a)^2$ for the raw
second moments). Estimating $\Theta$ therefore reduces to estimating linear and
bilinear functionals of the CATE pair $(\tau_S,\tau_Y)$.

\subsection{Assumption ledger, general case}

Assumptions~\ref{ass:iid}--\ref{ass:geometry} carry over verbatim as the general-case
readings (with $\Psi$ and $C$ in place of the correlation map $\varphi$ and the matrix
$\Sigma$). We state in full the two assumptions genuinely new relative to
Section~\ref{sec:layer1}.

Assumption~\ref{ass:A0} (A0, conditional-effect transportability given $X$) is
unchanged; it is what licenses \eqref{eq:delta-linear} on a continuous $\cX$ exactly
as on a finite one.

\begin{assumption}[A5: Bilinear-functional $\sqrt n$-estimability]\label{ass:elbow}
Each bilinear functional $\Theta_{ab}$ admits a $\sqrt n$-consistent, asymptotically
linear (debiased) estimator. Let $\tau_S,\tau_Y$ lie in H\"older/Sobolev balls of
smoothness $s_S,s_Y$ on $\cX\subseteq\R^d$, and suppose $C$ is identity-like
(non-smoothing) rather than finite rank or trace-class with decaying spectrum ---
Section~\ref{sec:rank} treats the case where $C$ is not identity-like, in which the
present assumption is not the relevant condition. For the identity (Dirac) kernel,
the established minimax boundary for $\sqrt n$-estimability of a quadratic/bilinear
functional applies moment-by-moment: the off-diagonal moment $\Theta_{SY}$ (a
bilinear functional of two \emph{different} functions) has boundary $s_S+s_Y>d/2$,
while each diagonal moment $\Theta_{aa}$ (a quadratic functional of a single
function) has the sharper boundary $s_a>d/4$, both attained by
higher-order-influence-function or rate-double-robust ("undersmoothed")
estimators specifically constructed to minimize the \emph{product} of the two CATE
estimation errors rather than each error's individual mean-squared risk
\citep{BickelRitov1988,BirgeMassart1995,Robins2008HOIF,Robins2017Minimax,Kennedy2024Minimax}.
Since $\hat\Theta=\Psi(\hat m)$ requires all three moments to be $\sqrt n$-estimable
simultaneously, the binding combined condition is
\begin{equation}\label{eq:elbow}
    \min(s_S,s_Y) > d/4,
\end{equation}
which implies the off-diagonal boundary automatically
($s_S+s_Y\ge2\min(s_S,s_Y)>d/2$) and so subsumes it; \eqref{eq:elbow} is
\emph{sufficient} for $\sqrt n$-estimability of $\Theta$ but is not shown here to be
\emph{necessary} for the ratio $\Theta$ itself (a shared scale factor in the three
moments can in principle cancel, as $a(\lambda)$ does in
Proposition~\ref{prop:interior}, so componentwise minimax lower bounds for the three
moments individually do not automatically transfer to a necessity statement for
$\Theta$). Below \eqref{eq:elbow}, the minimax rate for the moment(s) that fail it is
$n^{-4s/(4s+d)}$ ($s=\min(s_S,s_Y)$ in the relevant sense), strictly slower than
$n^{-1/2}$, and no $\sqrt n$-consistent asymptotically normal estimator of $\Theta$
exists; under balanced smoothness $s_S=s_Y=s$, \eqref{eq:elbow} reduces to $s>d/4$
exactly (not merely as a special case of a weaker general statement, since the
diagonal and off-diagonal boundaries coincide at balance). We adopt \eqref{eq:elbow}
as the working condition throughout, citing this literature for its attainability;
Lemma~\ref{lem:elbow-remainder} below proves only a \emph{weaker, naive} sufficient
condition using the simple one-step estimator with mean-squared-error-optimal CATE
tuning, and is explicit about not closing the resulting gap
(Remark~\ref{rem:elbow-gap}).
\end{assumption}

\begin{remark}[A5 is conservative; automatic on finite $X$]\label{rem:A5-conservative}
Condition~\eqref{eq:elbow} is the boundary for the Dirac kernel; a genuinely smooth
kernel $C$ typically \emph{relaxes} it (its integral operator regularizes the
representer $h_b=\int C\tau_b\,d\PP_0$ below), so \eqref{eq:elbow} is, at worst, a
conservative sufficient condition for our (bounded, geometry-induced) $C$. On a finite
$\cX$ every function is arbitrarily smooth and $\Theta_{ab}=\tau_a^\top\Sigma\tau_b$ is
a fixed polynomial of finitely many cell means, so \eqref{eq:elbow} is not the relevant
condition there --- $d$ is not even defined on a finite $\cX$. Section~\ref{sec:rank}
identifies the actual mechanism: $\Sigma$ has finite rank, and finite rank (not
finiteness of $\cX$ per se) is what removes the elbow, which is why
Section~\ref{sec:layer1} never states Assumption~\ref{ass:elbow}.
\end{remark}

\subsection{Effective rank of the reweighting-covariance kernel}\label{sec:rank}

Assumption~\ref{ass:elbow}'s elbow is a property of the kernel $C$, not of $\cX$ per
se; this section makes that precise and identifies the two endpoints between which
the general theory of this section and the finite-support theory of
Section~\ref{sec:layer1} sit as special cases of one underlying structure. Write the
spectral (Mercer) decomposition
\begin{equation}\label{eq:rank-decomp-kernel}
    C = \sum_j a_j\,\phi_j\otimes\phi_j
\end{equation}
on $L^2(\PP_0)$, i.e.\ $C$ is self-adjoint \emph{with respect to the $L^2(\PP_0)$
inner product} (on a finite $\cX$ this does not mean $C$'s matrix representation is
symmetric in the naive sense when $p_0$ is non-uniform, since the inner product
itself carries the $p_0$-weighting). Then, substituting into
\eqref{eq:reduction-bilinear},
\begin{equation}\label{eq:rank-decomp}
    \Theta_{ab} = \sum_j a_j\,\langle\tau_a,\phi_j\rangle\,\langle\tau_b,\phi_j\rangle.
\end{equation}
Two endpoints of \eqref{eq:rank-decomp} delimit the elbow's scope:
\begin{itemize}
    \item \textbf{Finite rank.} If $C$ has finite rank $r$ (only finitely many
    $a_j\ne0$), each $\Theta_{ab}$ is a finite sum of products of \emph{linear}
    functionals $\langle\tau_a,\phi_j\rangle$; each is estimable at the ordinary
    double-robust product nuisance rate $\|\hat\pi-\pi\|\cdot\|\hat m-m\|=o_P(n^{-1/2})$
    with no \emph{information-theoretic} elbow, for any $d$ --- finite rank removes
    the impossibility region of Assumption~\ref{ass:elbow}, but the ordinary
    nuisance-rate conditions for estimating each $\langle\tau_a,\phi_j\rangle$ still
    apply. Proposition~\ref{prop:interior}'s kernel $\Sigma=a(\lambda)(I_K-\mathbf1\mathbf1^\top/K)$
    is finite rank: $I_K-\mathbf1\mathbf1^\top/K$ is the orthogonal projection onto
    $\mathbf1^\perp$ (eigenvalue $0$ on $\mathrm{span}\{\mathbf1\}$, eigenvalue $1$
    with multiplicity $K-1$ on its complement), and viewed as an operator on
    $L^2(\PP_0)$, $C=a\cdot D^{-1}(I_K-\mathbf1\mathbf1^\top/K)$ with $D=\mathrm{diag}(p_0)$
    invertible, so the rank is $K-1$ regardless of $p_0$'s composition --- though the
    eigenbasis itself does depend on $p_0$ when $p_0$ is non-uniform (cell indicators
    are eigenfunctions of $C$ only when $p_0$ is uniform). It is this rank statement,
    basis-free, that recovers Proposition~\ref{prop:interior} as this section's
    finite-rank endpoint; finiteness of $\cX$ was never the operative reason
    Section~\ref{sec:layer1} needed no elbow.
    \item \textbf{Full rank / identity-like.} If $C$ has infinitely many nonzero
    $a_j$ not decaying to $0$ (the Dirac-kernel case Assumption~\ref{ass:elbow}
    addresses), $\Theta_{ab}$ is a genuinely quadratic/bilinear functional of
    $\tau_a,\tau_b$ and the elbow of Assumption~\ref{ass:elbow} is unavoidable.
\end{itemize}
A growing-truncation sieve at level $J_n$ interpolates the two endpoints, with tail
bias $\sum_{j>J_n}a_j\langle\tau_a,\phi_j\rangle\langle\tau_b,\phi_j\rangle=o_P(n^{-1/2})$
required for $\sqrt n$-estimability of $\Theta_{ab}$ (the diagonal tail terms $a=b$
are $\sum_{j>J_n}a_j\langle\tau_a,\phi_j\rangle^2\ge0$, non-negative with no
Cauchy--Schwarz relief, consistent with the diagonal moments binding in
Assumption~\ref{ass:elbow}). This document does not derive the rate at which $J_n$
must grow relative to $n$ to close this bias against the nuisance rates of
Section~\ref{sec:layer2}: that derivation, together with a full account of mixed
discrete/continuous covariates ($C=\Sigma\oplus C_{\mathrm{cont}}$, simultaneously
finite rank in the discrete directions and elbow-bound in the continuous ones) is
left to future work; \eqref{eq:rank-decomp} states the two endpoints and the
mechanism connecting them, not the interpolating rate.

\subsection{CATE asymptotic linearity}

\begin{lemma}[A1: CATE asymptotic linearity]\label{lem:cate-linearity}
For each $a\in\{S,Y\}$ and each $x\in\cX$, the CATE estimator admits the expansion
\begin{equation}\label{eq:cate-expansion}
    \hat\tau_a(x)-\tau_a(x) = \frac1n\sum_{i=1}^n\IF_{\tau_a}(O_i;x) + R_a(x),
\end{equation}
with $\E_{\PP_0}[\IF_{\tau_a}(O;x)]=0$ for every $x$ and
$\|R_a\|_{L^2(\PP_0)}=o_P(n^{-1/4})$. For a randomized trial $\IF_{\tau_a}$ is the
per-$x$ difference-of-means score; for AIPW/DR-learners it is the pointwise AIPW
score. On a finite $\cX$, $\IF_{\tau_a}(O;k)$ is exactly the per-cell AIPW score of
\eqref{eq:aipw-estimator}, integrated against a point mass at $k$: this is the bridge
to Section~\ref{sec:layer1}.
\end{lemma}

\begin{remark}[A concrete kernel form of $\IF_{\tau_a}$ on continuous $\cX$]\label{rem:kernel-IF-concrete}
Lemma~\ref{lem:cate-linearity} asserts that, for AIPW/DR-learners, $\IF_{\tau_a}$ is
``the pointwise AIPW score'' without specifying an estimator. We make this concrete
for continuous $\cX\subseteq\R^d$, generalizing the finite-$\cX$ per-cell estimator of
Lemma~\ref{lem:cancellation} (equation~\eqref{eq:aipw-estimator}), whose per-unit
bracket is named $\mathrm{score}_S(O)$ in the proof of Theorem~\ref{thm:main}
(Appendix~\ref{app:proof-main}, eq.~\eqref{eq:main-proof-expansion}). For outcome
type $a\in\{S,Y\}$ let $\mathrm{score}_a(O)$ denote that same cross-fitted AIPW
bracket for $\tau_a$ (matching $\mathrm{score}_S(O)$ when $a=S$; defined analogously
with $\mu_1^Y,\mu_0^Y$ and $Y$ when $a=Y$). Estimate $\tau_a(x)$ by kernel (local
polynomial) regression of $\mathrm{score}_a(O)$ on $X$, using a kernel $K$ of order
$\ell\geq\lceil s_a\rceil$ (equivalently, local polynomial degree
$\geq\lceil s_a\rceil-1$) so that the estimator's bias matches the Hölder-$s_a$
smoothness of $\tau_a$; rescaled kernel $K_h(u)=h^{-d}K(u/h)$, bandwidth $h=h_n\to0$,
$nh^d\to\infty$. The pointwise influence function generalizes the
finite-$\cX$ weight $\mathbf 1\{X{=}k\}/p_0(k)$ of \eqref{eq:main-proof-expansion} to
a kernel weight:
\begin{equation}\label{eq:kernel-IF}
    \IF_{\tau_a}(O_i;x) = \frac{K_h(X_i-x)}{f(x)}\Big[\mathrm{score}_a(O_i)-\tau_a(x)\Big],
\end{equation}
where $f$ is the covariate density. Under standard nonparametric regularity---$f$
bounded away from $0$ and $\infty$ on $\cX$, $\tau_a$ H\"older-$s_a$, and
$\Var[\mathrm{score}_a(O)\mid X{=}x]$ bounded---the classical bias--variance
decomposition for kernel/local-linear regression \citep{Stone1980}, specialized to
CATE estimation in \citet{Kennedy2024Minimax}, gives
\begin{equation}\label{eq:kernel-bias-variance}
    \hat\tau_a(x)-\tau_a(x) = \underbrace{\frac1n\sum_{i=1}^n\IF_{\tau_a}(O_i;x)}_{O_P\left((nh^d)^{-1/2}\right)}
        + \underbrace{O(h^{s_a})}_{\text{bias}}.
\end{equation}
The mean-squared-error-optimal bandwidth $h\asymp n^{-1/(2s_a+d)}$ balances both
terms at rate $n^{-s_a/(2s_a+d)}$, exactly the rate Lemma~\ref{lem:elbow-remainder}
already assumes for $\|\hat\tau_a-\tau_a\|_{L^2(\PP_0)}$: this construction is what
makes that rate achievable rather than an unjustified black box.

The same bandwidth choice also settles Lemma~\ref{lem:cate-linearity}'s own remainder
condition, rather than leaving it as an independent assumption. At
$h\asymp n^{-1/(2s_a+d)}$, $R_a(x)$ in \eqref{eq:cate-expansion} is exactly the bias
term of \eqref{eq:kernel-bias-variance}, and
\[
    n^{-s_a/(2s_a+d)} = o\big(n^{-1/4}\big)
    \iff \frac{s_a}{2s_a+d} > \frac14
    \iff 4s_a>2s_a+d
    \iff s_a>\frac d2.
\]
Under balanced smoothness $s_S=s_Y=s$, $s>d/2$ is exactly $s_S+s_Y>d$: the naive
elbow of Assumption~\ref{ass:elbow} and Lemma~\ref{lem:elbow-remainder}. So
$\|R_a\|_{L^2(\PP_0)}=o_P(n^{-1/4})$ is not a separate regularity condition here---it
holds if and only if the naive elbow condition used everywhere else in this document
holds, at the MSE-optimal bandwidth.

This construction attains only the naive elbow $s_S+s_Y>d$, not the sharp boundary
$s_S+s_Y>d/2$ of Assumption~\ref{ass:elbow}: reaching the sharp boundary requires the
higher-order-influence-function or undersmoothed constructions cited there, which we
do not construct here (Remark~\ref{rem:elbow-gap}).
\end{remark}

\subsection{The bilinear-functional efficient influence function}

\begin{lemma}[A3: EIF of a bilinear functional of CATEs]\label{lem:bilinear-eif}
Fix the kernel $C$ and the covariate law of $\PP_0$ (Convention~\ref{conv:geometry}).
Suppose $\tau_a,\tau_b$ have influence functions $\IF_{\tau_a},\IF_{\tau_b}$
(Lemma~\ref{lem:cate-linearity}) and $\E_{\PP_0}[\tau_a^2],\E_{\PP_0}[\tau_b^2]<\infty$.
Then $\Theta_{ab}=\iint C(x,x')\tau_a(x)\tau_b(x')\,d\PP_0(x)d\PP_0(x')$ is pathwise
differentiable with mean-zero efficient influence function
\begin{equation}\label{eq:chi-ab}
    \chi_{ab}(O) = \int h_b(x)\,\IF_{\tau_a}(O;x)\,d\PP_0(x)
        + \int h_a(x')\,\IF_{\tau_b}(O;x')\,d\PP_0(x'),
\end{equation}
where
\begin{equation}\label{eq:h-directions}
    h_b(x)=\int C(x,x')\tau_b(x')\,d\PP_0(x'), \qquad h_a(x')=\int C(x',x)\tau_a(x)\,d\PP_0(x)
\end{equation}
are the two \emph{directional representers}. When $a=b$,
$\chi_{aa}(O)=2\int h_a(x)\IF_{\tau_a}(O;x)\,d\PP_0(x)$ (the factor $2$ reflecting that
$\Theta_{aa}$ is a pure quadratic). On a finite $\cX$, $h_b(k)=(\Sigma\tau_b)_k/p_0(k)$
and \eqref{eq:chi-ab} reduces exactly to
\begin{equation}\label{eq:chi-ab-discrete}
    \chi_{ab}(O) = (\Sigma\tau_b)^\top\mathbf{IF}_{\tau_a}(O) + (\Sigma\tau_a)^\top\mathbf{IF}_{\tau_b}(O),
    \qquad \mathbf{IF}_{\tau_a}(O)=(\IF_{\tau_a}(O;1),\dots,\IF_{\tau_a}(O;K))^\top,
\end{equation}
verified directly in Section~\ref{sec:recovery}.
\end{lemma}

The proof (Appendix~\ref{app:proof-bilinear}) computes the pathwise derivative of
$\Theta_{ab}$ along a one-dimensional submodel by the product rule over the two
$\tau$-factors (Leibniz, justified by bounded $C$ and square-integrable $\tau$), then
identifies the Riesz representer of the resulting linear functional of the submodel's
score; symmetrizing over the two factors gives \eqref{eq:chi-ab}, and specializing the
integrals to point masses on a finite $\cX$ gives \eqref{eq:chi-ab-discrete}
(Appendix~\ref{app:proof-recovery} verifies this specialization by direct
substitution).

\subsection{One-step debiased estimator and the estimability elbow}

\begin{lemma}[A4: debiased cross-fit estimator]\label{lem:onestep}
With cross-fitted $\hat\tau_a,\hat\tau_b,\hat C$ and
$\hat h_b(x)=\int\hat C(x,x')\hat\tau_b(x')\,d\PP_0(x')$, the one-step estimator
\begin{equation}\label{eq:onestep}
    \hat\Theta_{ab} = \iint \hat C(x,x')\hat\tau_a(x)\hat\tau_b(x')\,d\PP_0(x)d\PP_0(x')
        + \frac1n\sum_{i=1}^n\hat\chi_{ab}(O_i)
\end{equation}
is Neyman-orthogonal in the CATE pair $\beta=(\tau_a,\tau_b)$ (the correction cancels
the plug-in's first-order bias, by construction of $\chi_{ab}$ as its canonical
gradient) and, because the representers $h_a,h_b$ are functions of $x$ alone that do
not depend on the outcome-model nuisances $\eta=(e,\mu_0,\mu_1)$, inherits Neyman
orthogonality in $\eta$ as well, by the same ``orthogonality survives reweighting''
argument used for Theorem~\ref{thm:main}'s AIPW score (Appendix~\ref{app:proof-onestep}).
\end{lemma}

\begin{lemma}[A5: bilinear remainder and the elbow]\label{lem:elbow-remainder}
The estimator \eqref{eq:onestep} satisfies
$\hat\Theta_{ab}-\Theta_{ab}=\frac1n\sum_i\chi_{ab}(O_i)+R_n$ up to $o_P(n^{-1/2})$
nuisance terms, where the bilinear-functional remainder is exactly
\begin{equation}\label{eq:remainder}
    R_n = \iint \hat C(x,x')\,(\hat\tau_a-\tau_a)(x)\,(\hat\tau_b-\tau_b)(x')\,d\PP_0(x)d\PP_0(x'),
\end{equation}
with no higher-order terms (a bilinear functional's Taylor expansion terminates
exactly at the product term). By Cauchy--Schwarz and the operator-norm bound
$\|\hat C\|_{op}\le C_w^2$ (Assumptions~\ref{ass:moments},\ref{ass:geometry}),
\begin{equation}\label{eq:CS-remainder}
    |R_n| \le \|\hat C\|_{op}\,\|\hat\tau_a-\tau_a\|_{L^2(\PP_0)}\,\|\hat\tau_b-\tau_b\|_{L^2(\PP_0)},
\end{equation}
so $R_n=o_P(n^{-1/2})$ whenever the product of CATE $L^2$-rates is $o_P(n^{-1/2})$.
Using mean-squared-error-optimal nonparametric CATE estimators, each
$\|\hat\tau_a-\tau_a\|_{L^2(\PP_0)}=O_P(n^{-s_a/(2s_a+d)})$. For the diagonal moments
($a=b$) the product is $o_P(n^{-1/2})$ exactly under $s_a>d/2$; for the off-diagonal
moment ($a\ne b$) exactly under $s_S/(2s_S+d)+s_Y/(2s_Y+d)>1/2$, which reduces to the
arithmetic condition $s_S+s_Y>d$ only under balanced smoothness $s_S=s_Y$ and is a
strictly weaker (less restrictive) requirement than $s_S+s_Y>d$ under imbalance.
The single combined condition
\begin{equation}\label{eq:elbow-naive}
    \min(s_S,s_Y) > d/2
\end{equation}
(Remark~\ref{rem:elbow-gap}) is sufficient for all three moments simultaneously: it
delivers the diagonal condition directly, and since $s\mapsto s/(2s+d)$ is
increasing, $s_a>d/2$ for each $a\in\{S,Y\}$ gives $s_a/(2s_a+d)>1/4$ for each, so
their sum exceeds $1/2$, delivering the off-diagonal condition as well.
\eqref{eq:elbow-naive} is exactly twice Assumption~\ref{ass:elbow}'s stated combined
boundary $\min(s_S,s_Y)>d/4$.
\end{lemma}

\begin{remark}[The naive bound proved here does not reach the literature's elbow]\label{rem:elbow-gap}
The Cauchy--Schwarz argument above, using each $\hat\tau_a$ tuned to its own
mean-squared-error-optimal nonparametric rate, gives $R_n=o_P(n^{-1/2})$ for the
diagonal moments only under $s_a>d/2$ (each $a$ separately) and for the
off-diagonal moment only under $s_S/(2s_S+d)+s_Y/(2s_Y+d)>1/2$ --- not the sharper
diagonal boundary $s_a>d/4$ or off-diagonal boundary $s_S+s_Y>d/2$ stated as
Assumption~\ref{ass:elbow}. The single condition $\min(s_S,s_Y)>d/2$
\eqref{eq:elbow-naive} is sufficient for all three simultaneously and is exactly
twice Assumption~\ref{ass:elbow}'s combined threshold $\min(s_S,s_Y)>d/4$; under
balanced smoothness $s_S=s_Y=s$ this is the naive threshold $s>d/2$ versus the sharp
threshold $s>d/4$, matching the single-moment gap already familiar from the
finite-support special case. Reaching the sharp boundary requires replacing the
plug-in one-step estimator with a higher-order-influence-function or undersmoothed
construction that optimizes the \emph{product} $\|\delta_a\|\|\delta_b\|$ directly
rather than each factor's individual risk --- exactly the machinery of
\citet{Robins2008HOIF,Robins2017Minimax} cited in Assumption~\ref{ass:elbow}, which
we do not reproduce here. Consequently, Theorem~\ref{thm:general} as proved in this
document (via Lemma~\ref{lem:elbow-remainder}'s naive bound) is rigorously
established for $\min(s_S,s_Y)>d/2$; extending the rigorous range down to the stated
elbow $\min(s_S,s_Y)>d/4$ is a real, currently unclosed step in this document,
resolved in the cited literature by a different (higher-order) estimator than the
one-step construction of Section~\ref{sec:layer2}, and should be read as such rather
than as fully re-derived here.
\end{remark}

\subsection{Moment-vector influence function and delta-method composition}

\begin{lemma}[A7: moment-vector EIF]\label{lem:moment-eif}
Stacking the linear-moment scores
$\psi_{\E_\mu\Delta_a}(O)=\int\bar a(x)\IF_{\tau_a}(O;x)\,d\PP_0(x)$, the bilinear
scores $\chi_{SS},\chi_{YY},\chi_{SY}$ of Lemma~\ref{lem:bilinear-eif}, and the
chain-rule score for raw second moments gives an influence function $\psi_m$ for the
moment vector $m$, and, by the functional delta method under Hadamard
differentiability of $\Psi$ with bounded gradient on
$\{\Var_\mu\Delta_S,\Var_\mu\Delta_Y\ge c_0\}$ (Assumption~\ref{ass:nondegeneracy}),
\begin{equation}\label{eq:psi-theta-general}
    \psi_\Theta(O) = \nabla\Psi(m)^\top\psi_m(O).
\end{equation}
Because $\chi_{SY}$ pairs $\IF_{\tau_S}$ and $\IF_{\tau_Y}$ evaluated at the
\emph{same} unit $O$, $\Var(\psi_\Theta)$ is not block-diagonal across $S$ and $Y$.
\end{lemma}

\subsection{Main result}

\begin{theorem}[General asymptotic normality]\label{thm:general}
Under A0, A5 (Assumptions~\ref{ass:A0},\ref{ass:elbow}) and the general-case readings
of Assumptions~\ref{ass:iid}--\ref{ass:geometry}, with one-step cross-fit debiased
moment estimators (Lemmas~\ref{lem:cate-linearity}--\ref{lem:elbow-remainder}) and
$\hat\Theta=\Psi(\hat m)$, conditional on the sampled geometry, as $n\to\infty$,
\begin{equation}\label{eq:thm-general}
    \sqrt n\big(\hat\Theta-\Theta(\Qhat)\big)
    = \frac1{\sqrt n}\sum_{i=1}^n\psi_\Theta(O_i;\Qhat) + O_P\big((n/M)^{1/2}\big) + o_P(1)
    \xrightarrow{d} N\big(0,\E_{\PP_0}[\psi_\Theta^2]\big),
\end{equation}
with $\psi_\Theta$ of \eqref{eq:psi-theta-general}. If additionally $n=o(M)$, the
Monte Carlo term is $o_P(1)$ and unconditionally
$\sqrt n(\hat\Theta-\Theta(\mu))\xrightarrow{d}N(0,\sigma^2(\mu))$.
\end{theorem}

The proof (Appendix~\ref{app:proof-general}) combines the moment expansions of
Lemmas~\ref{lem:cate-linearity} and~\ref{lem:elbow-remainder}, the delta method of
Lemma~\ref{lem:moment-eif} (with the ratio-type $\Psi$ for correlation closing under a
bounded Hessian on the non-degenerate region, Assumption~\ref{ass:nondegeneracy}), the
i.i.d.\ CLT conditional on the geometry, and the Markov-chain CLT for the Monte Carlo
term (Lemma~\ref{lem:mcmc-error}) --- the same four-step structure as
Theorem~\ref{thm:main}, generalized from the matrix $\Sigma$ to the operator $C$.

% ===========================================================================
\section{Recovery: the finite-support theorem as a special case}\label{sec:recovery}
% ===========================================================================

\begin{corollary}[Recovery of Theorem~\ref{thm:main}]\label{cor:recovery}
When $\cX=\{1,\dots,K\}$ is finite, $C\to\Sigma$ as in \eqref{eq:sigma-def},
$\IF_{\tau_a}$ is the per-cell AIPW score of \eqref{eq:aipw-estimator}, and $\Psi$ is
the correlation map $\varphi$ of Section~\ref{sec:layer1}, Theorem~\ref{thm:general}
reduces exactly to Theorem~\ref{thm:main}, with the same influence function
$\psi_\Theta$ and the same $\sigma^2(\Qhat)$.
\end{corollary}

\begin{proof}[Proof sketch; full derivation in Appendix~\ref{app:proof-recovery}]
On finite $X$, A5 is automatic (Remark~\ref{rem:A5-conservative}) and
$\Theta_{ab}=\tau_a^\top\Sigma\tau_b$ by \eqref{eq:discrete-quadform}. Substituting the
discrete kernel into \eqref{eq:h-directions}, $h_b(k)=(\Sigma\tau_b)_k/p_0(k)$, and
into \eqref{eq:chi-ab}, using $\int g(x)\IF_{\tau_a}(O;x)\,d\PP_0(x)=\sum_kp_0(k)g(k)\IF_{\tau_a}(O;k)$
for any bounded $g$, gives exactly \eqref{eq:chi-ab-discrete}
(Lemma~\ref{lem:recovery-chi}, Appendix~\ref{app:proof-recovery}): the $p_0(k)$
introduced by the representer's denominator and the $p_0(k)$ from the integral
measure cancel, leaving the pure matrix expression $(\Sigma\tau_b)^\top\mathbf{IF}_{\tau_a}(O)+(\Sigma\tau_a)^\top\mathbf{IF}_{\tau_b}(O)$.
Composing with $\varphi$'s gradient (the chain-rule composition of $\nabla\varphi(m(\mu))$
with the moment derivatives, exactly as used in \eqref{eq:psi-theta-layer1}) and substituting into \eqref{eq:psi-theta-general}
reproduces \eqref{eq:psi-theta-layer1}. All remainders are those of
Lemmas~\ref{lem:cate-linearity} and~\ref{lem:elbow-remainder}, specialized to the
finite case, so the two limit laws coincide.
\end{proof}

\begin{corollary}[Structural correspondence with the plug-in debiasing correction]\label{cor:onestep-match}
The plug-in estimator $\hat\Theta_{ab}^{\mathrm{plugin}}=\hat\tau_a^\top\Sigma\hat\tau_b$
(no correction) has leading bias
$\E[\hat\Theta_{ab}^{\mathrm{plugin}}-\Theta_{ab}]=\tr(\Sigma V)/n+O(n^{-3/2})$, where
$V_{kk'}=\E_{\PP_0}[\IF_{\tau_a}(O;k)\IF_{\tau_b}(O;k')]$ is the CATE-score
covariance (Appendix~\ref{app:proof-recovery}, Lemma~\ref{lem:plugin-bias}). The
package's optional one-step correction (\texttt{debiased = TRUE} in
\texttt{tv\_ball\_correlation\_IF\_adaptive()}) estimates $\tr(\Sigma V)/n$ for the
covariance term $\Theta_{SY}$ and propagates it to the correlation scale via
$\partial\rho/\partial\Theta_{SY}=1/(\mathrm{sd}_S\,\mathrm{sd}_Y)$, adding the result
to the plug-in correlation. This targets the same $O(1/n)$ bias term identified by
Lemma~\ref{lem:plugin-bias} and is the same class of correction as the general
one-step estimator \eqref{eq:onestep} applied at the finite-support recovery of
Corollary~\ref{cor:recovery}: both use the bilinear functional's exact second-order
structure (Lemma~\ref{lem:elbow-remainder}) to target the plug-in's $O(1/n)$ bias
rather than merely bound it. Remark~\ref{rem:onestep-scope} states what this
structural correspondence does, and does not, establish.
\end{corollary}

\begin{remark}[Scope of the structural correspondence]\label{rem:onestep-scope}
Corollary~\ref{cor:onestep-match} asserts the two corrections are the \emph{same
structural object}; it does not assert they are numerically identical estimators
once both are fully specified in finite samples, for two distinct reasons. First,
\eqref{eq:onestep}'s one-step term $\frac1n\sum_i\hat\chi_{ab}(O_i)$, when specialized
via Lemma~\ref{lem:recovery-chi} and expanded using the per-cell asymptotic linearity
of $\hat\tau_a(k)-\tau_a(k)\approx\frac1n\sum_i\IF_{\tau_a}(O_i;k)$
(Lemma~\ref{lem:cate-linearity}), involves the same $n_k$-vs-$n$ per-cell
renormalization bookkeeping that produced Lemma~\ref{lem:cancellation}'s exact
cancellation in Theorem~\ref{thm:main}'s estimator. The two corrections share the
same leading-order bias \emph{magnitude} ($\tr(\Sigma V)/n$ in both cases,
Lemma~\ref{lem:plugin-bias}), but a term-for-term match through this renormalization
is a separate, unclosed step. Second, and more concretely: propagating the
covariance bias $\tr(\Sigma V)/n$ to the correlation scale via
$\partial\rho/\partial\Theta_{SY}$ alone does not separately correct the
denominator moments $\Theta_{SS},\Theta_{YY}$ (each of which has its own
$\tr(\Sigma V_{aa})/n$ bias by Lemma~\ref{lem:plugin-bias}, applied with $a=b$); a
full correlation-scale correction would propagate all three moments' biases through
$\rho=\Theta_{SY}/\sqrt{\Theta_{SS}\Theta_{YY}}$'s gradient, not $\Theta_{SY}$'s bias
alone. Corollary~\ref{cor:onestep-match} should therefore be read as establishing
structural correspondence between two constructions that target the same $O(1/n)$
bias term, not as certifying an exact, fully correlation-scale-propagated
finite-sample equivalence between them.
\end{remark}

% ===========================================================================
\section{Discussion}\label{sec:discussion}
% ===========================================================================

\paragraph{What the estimand measures, and what it assumes.} $\Theta(\PP_0,\lambda)$
answers a different question than the proportion of treatment effect (PTE) or
within-study mediation: it asks whether treatment effects on the surrogate and the
outcome co-vary across plausible future population compositions, not how much of the
effect in the current study is mediated through the surrogate. Assumption~\ref{ass:A0}
scopes this to compositional shifts in measured covariates; under the interpretive
reading of Remark~\ref{rem:A0-reading}, it is an untestable-from-one-study assumption
that unmeasured study-level effect modifiers do not exist, and $\lambda$-sensitivity
analysis (reporting $\hat\Theta_n(\lambda)$ across a grid) is the practical proxy for
robustness to violations, not a substitute for the assumption itself.

\paragraph{Three regimes, three interpretations of $\lambda$.} Propositions~\ref{prop:interior}
and~\ref{prop:saturated}, together with Theorem~\ref{thm:main}, answer ``what does
$\lambda$ do?'' concretely: for $\lambda$ small enough that no covariate stratum can
be driven to zero mass (\emph{interior}), and again for $\lambda$ large enough that
the TV ball covers the whole simplex (\emph{saturated}), $\Theta$ is exactly the
stratum-level meta-analytic correlation, independent of both $\lambda$ and $\PP_0$'s
composition; only in the \emph{intermediate} range does the simplex boundary make it
asymmetrically harder to reweight rare strata upward without exhausting common ones,
so that $\Theta$ need not equal this closed form and may genuinely depend on how the
observed population is distributed across strata. The paper's validation study
operates in the intermediate regime (Remark~\ref{rem:interior-scope}), and
Example~\ref{ex:binding} confirms the general, boundary-aware machinery of
Theorem~\ref{thm:main} does genuine work there, not merely recovering the closed
form. For continuous $X$, the interior regime as defined does not exist
(Lemma~\ref{lem:interior-atomic}), and the saturated-regime closed form
(Proposition~\ref{prop:saturated}) does not directly generalize either, since it is
stated and proved only for finite $K$ and the uniform measure $\mu$ needed to state
a continuous analog is not well-defined on an infinite-dimensional simplex
(Remark~\ref{rem:no-continuous-interior}); Proposition~\ref{prop:sieve-limit} gives
the corresponding continuous-$X$ limit via an equal-$\PP_0$-mass sieve instead.

\paragraph{What is conditional, and why.} Every result in this document conditions on
the sampled geometry (Convention~\ref{conv:geometry}), which subsumes both the
Monte Carlo approximation (finite $M$) and the use of an estimated center $\Pn$ rather
than $\PP_0$ (Remark~\ref{rem:geometry-scope}). The unconditional statements
(Theorem~\ref{thm:main}(b), Theorem~\ref{thm:general}'s closing sentence) require
$M$ to grow strictly faster than $n$; absent that, reported inference targets the
$M$-study, $\Pn$-centered correlation, not the population quantity, and should be
labeled as such (Remark~\ref{rem:practical-use}).

\paragraph{Limitations genuinely left open.} (i) A fully unconditional theory
differentiating the kernel $C$'s and the geometry's own dependence on $\PP_0$
(Remarks~\ref{rem:geometry-scope},\ref{rem:kernel-dependence}) is not derived here.
(ii) The exact finite-sample correspondence of Corollary~\ref{cor:onestep-match} is
flagged, not closed (Remark~\ref{rem:onestep-scope}). (iii) Below the estimability
elbow \eqref{eq:elbow}, no $\sqrt n$-consistent estimator of the general estimand
exists; higher-order influence functions attain the minimax rate but do not restore
asymptotic normality, and undersmoothed, conservative intervals are the appropriate
fallback. (iv) Extending the absolute-continuity restriction $\Q\ll\PP_0$ to permit
genuine extrapolation to unobserved covariate values requires additional modeling
assumptions not developed here.

% ===========================================================================
\appendix
\renewcommand{\thesection}{\Alph{section}}
\section{Auxiliary results}\label{app:aux}
% ===========================================================================

\begin{lemma}[Norm equivalence]\label{lem:normequiv}
Under Assumption~\ref{ass:moments}, $\|w\|_{L^2(\PP_0)}\le C_w$ and, for any bounded
$g$ with $\|g\|_\infty\le B$, $\|w\cdot g\|_{L^2(\PP_0)}\le C_w B$.
\end{lemma}

\begin{lemma}[MCMC approximation error]\label{lem:mcmc-error}
Under Assumption~\ref{ass:geometry}, for any bounded measurable
$g:\mathrm{supp}(\mu)\to\R$ with $\|g\|_\infty\le1$,
\[
    \big|\E_{\Qhat}[g(\Q)]-\E_\mu[g(\Q)]\big| = O_P(M^{-1/2}),
\]
via the Markov chain central limit theorem for the geometrically ergodic hit-and-run
sampler (not the i.i.d.\ CLT, since the $\Q_m$ are dependent draws): the asymptotic
variance is inflated by the chain's integrated autocorrelation time, but the rate is
unchanged from the independent case.
\end{lemma}

\begin{lemma}[Mean-zero AIPW influence function]\label{lem:if-meanzero}
For fixed $\Q_m$ and $\psi_S^{\mathrm{AIPW}}$ of \eqref{eq:psi-aipw},
$\E_{\PP_0}[\psi_S^{\mathrm{AIPW}}(O;\Q_m)]=0$ under either correct $e$ or correct
$(\mu_1^S,\mu_0^S)$ (double robustness), and $\partial_\eta\E_{\PP_0}[\psi_S^{\mathrm{AIPW}}]=0$
at the truth (Neyman orthogonality), surviving multiplication by the $\eta$-free
weight $w(X;\Q_m)$.
\end{lemma}

% ===========================================================================
\section{Proofs}\label{app:proofs}
% ===========================================================================

\subsection{Proof of Proposition~\ref{prop:interior}}\label{app:proof-interior}

\begin{proof}
In the interior regime (Definition~\ref{def:interior}), the feasible set for the
signed measure $V$ is exactly $\cB=\{V\in\R^K:\mathbf1^\top V=0,\|V\|_1\le2\lambda\}$,
independent of $\PP_0$; ``uniform on $\cB$'' means uniform with respect to
$(K{-}1)$-dimensional Lebesgue measure on the hyperplane $\{\mathbf1^\top V=0\}$
containing $\cB$ (the natural reference measure, since $\cB$ has empty interior, hence
Lebesgue measure zero, in $\R^K$). This set is invariant under every permutation
$\pi$ of $\{1,\dots,K\}$: both $\mathbf1^\top V=0$ and $\|V\|_1\le2\lambda$ are
permutation-symmetric functions of $V$'s coordinates. The uniform measure on any
permutation-invariant convex body (with respect to the permutation-invariant
Lebesgue measure on its ambient hyperplane) is itself permutation-invariant (the
volume element is unchanged by relabeling coordinates), so if $V\sim\mathrm{Uniform}(\cB)$
then $V\stackrel{d}{=}\Pi V$ for every permutation matrix $\Pi$. In particular
$\E[V]=\E[\Pi V]=\Pi\E[V]$ for every $\Pi$, and the only vector fixed by every
coordinate permutation is a multiple of $\mathbf1$; combined with $\mathbf1^\top
\E[V]=\E[\mathbf1^\top V]=0$, this forces $\E[V]=0$. Let $\Sigma=\Cov(V)=\E[VV^\top]$
(using $\E[V]=0$). Permutation invariance gives $\Sigma=\Pi\Sigma\Pi^\top$ for every
permutation matrix $\Pi$; the only matrices commuting with the full symmetric group's
representation are $\Sigma=aI_K+b\,\mathbf1\mathbf1^\top$ for scalars $a,b$ (Schur's
lemma for the permutation representation, or direct verification: any $\Sigma$
satisfying $\Sigma_{kk}=\Sigma_{ll}$ for all $k,l$ and $\Sigma_{kk'}=\Sigma_{ll'}$ for
all $k\ne k',l\ne l'$ has this form). Since $\mathbf1^\top V=0$ almost surely,
$\Sigma\mathbf1=\Cov(V,\mathbf1^\top V)=0$, which \emph{determines} $b$ in terms of
$a$ (rather than eliminating a $b\mathbf1\mathbf1^\top$ term outright): $(a+bK)\mathbf1=0$, i.e.\
$b=-a/K$. Hence
\[
    \Sigma = a\Big(I_K-\tfrac1K\mathbf1\mathbf1^\top\Big)
\]
for some scalar $a=a(\lambda)>0$ (positive since $V$ is non-degenerate for
$\lambda>0$). Substituting into $\Theta_{ab}=\tau_a^\top\Sigma\tau_b$
\eqref{eq:discrete-quadform},
\[
    \Theta_{ab} = a\Big(\tau_a^\top\tau_b - \tfrac1K(\mathbf1^\top\tau_a)(\mathbf1^\top\tau_b)\Big)
    = aK\Big(\tfrac1K\textstyle\sum_k\tau_a(k)\tau_b(k) - \bar\tau_a\bar\tau_b\Big)
    = aK\cdot\mathrm{Cov}_K(\tau_a,\tau_b),
\]
the scalar $aK$ times the unweighted empirical covariance of the $K$ cell CATEs.
Forming the ratio $\Theta_{SY}/\sqrt{\Theta_{SS}\Theta_{YY}}$ cancels $aK$ exactly,
giving \eqref{eq:interior-reduction}; $p_0(k)$ never entered the computation once
$\Sigma$'s form was fixed by symmetry, so the result does not depend on $\PP_0$'s
composition (only on the support, since cells with $p_0(k)=0$ carry no CATE and are
excluded from the geometry entirely).
\end{proof}

\subsection{Proofs of Lemma~\ref{lem:root-polytope} and
Proposition~\ref{prop:exit-interior}}\label{app:proof-exit-interior}

\begin{proof}[Proof of Lemma~\ref{lem:root-polytope}]
\emph{(a).} For $V\ne0$ in $H$ (so $\mathbf1^\top V=0$; $V=0$ satisfies every facet
inequality trivially and is the midpoint of any antipodal vertex pair), let
$S^+=\{k:V_k\ge0\}$, a proper nonempty subset since $V\ne0$ and $\sum_kV_k=0$ rule
out $S^+\in\{\emptyset,[K]\}$. Since $\sum_kV_k=0$, $\sum_{k\in
S^+}V_k=-\sum_{k\notin S^+}V_k=\sum_{k\notin S^+}(-V_k)$, and both sides equal
$\frac12\|V\|_1$ because $\|V\|_1=\sum_{k\in S^+}V_k+\sum_{k\notin S^+}(-V_k)$. For
any other $\emptyset\ne S\subsetneq[K]$, $\sum_{k\in S}V_k\le\sum_{k\in S^+}V_k$:
adding an index with $V_k<0$ or removing one with $V_k\ge0$ only decreases the sum.
Hence $\max_S\sum_{k\in S}V_k=\frac12\|V\|_1$, and $\|V\|_1\le2\lambda\iff\sum_{k\in
S}V_k\le\lambda$ for every proper nonempty $S$.

\emph{(b).} \emph{Spanning.} Fix $V\in\cB_\lambda$, $V\ne0$, and let
$T=\frac12\|V\|_1\in(0,\lambda]$. The vectors $(V_k^+)_{k\in P}$ and
$(V_k^-)_{k\in N}$, $P=\{k:V_k>0\}$, $N=\{k:V_k<0\}$, are two nonempty collections
of positive numbers with equal total mass $T$ (by the computation in (a)); the
independence coupling $\pi_{ij}=V_i^+V_j^-/T\ge0$ ($i\in P,j\in N$) satisfies
$\sum_j\pi_{ij}=V_i^+$, $\sum_i\pi_{ij}=V_j^-$, and
$V=\sum_{i\in P,j\in N}\pi_{ij}(e_i-e_j)=\lambda\sum_{i,j}(\pi_{ij}/\lambda)(e_i-e_j)$
with $\sum_{i,j}\pi_{ij}/\lambda=T/\lambda\le1$; padding the deficiency
$1-T/\lambda$ with any antipodal pair $\{\lambda(e_i-e_j),\lambda(e_j-e_i)\}$ in
equal weight (contributing $0$, since $\frac12[\lambda(e_i-e_j)+\lambda(e_j-e_i)]=0$)
exhibits $V$ as a convex combination of the $K(K-1)$ points
$\{\lambda(e_i-e_j)\}$; $V=0$ is itself such a combination. \emph{Extremality.}
Suppose $\lambda(e_i-e_j)=\frac12(u+w)$ for $u,w\in\cB_\lambda$. Since
$u,w\in\cB_\lambda$, $\|u\|_1\le2\lambda$ and $\|w\|_1\le2\lambda$, so
$\|u\|_1+\|w\|_1\le4\lambda$; since also $u+w=2\lambda(e_i-e_j)$, the triangle
inequality gives $\|u\|_1+\|w\|_1\ge\|u+w\|_1=2\lambda\|e_i-e_j\|_1=4\lambda$. Both
bounds force $\|u\|_1+\|w\|_1=4\lambda=\|u+w\|_1$: equality in the triangle
inequality, which for $\|\cdot\|_1$ holds coordinatewise
($|u_k|+|w_k|=|u_k+w_k|$ for every $k$) and forces $\|u\|_1=\|w\|_1=2\lambda$. For
$k\notin\{i,j\}$, $u_k+w_k=0$ together with $|u_k|+|w_k|=|u_k+w_k|=0$ gives
$u_k=w_k=0$; the zero-sum constraint then forces $u=(u_i,-u_i)$, $w=(w_i,-w_i)$ on
$\{i,j\}$ (listing only the two nonzero coordinates), so $\|u\|_1=2|u_i|=2\lambda$
and likewise $\|w\|_1=2|w_i|=2\lambda$ give $u_i,w_i\in\{-\lambda,+\lambda\}$, and
$u_i+w_i=2\lambda$ (the $i$-th coordinate of $u+w=2\lambda(e_i-e_j)$) forces
$u_i=w_i=\lambda$, the only sign choice compatible with the sum: $u=w=\lambda(e_i-e_j)$.
By the spanning claim these $K(K-1)$ points are all the extreme points of
$\cB_\lambda$.

\emph{(c).} Among the vertices $\{\lambda(e_i-e_j)\}_{i\ne j}$, the $k$-th
coordinate equals $-\lambda$ when $j=k$ (the $K-1$ vertices $\lambda(e_i-e_k)$,
$i\ne k$), $+\lambda$ when $i=k$, and $0$ otherwise; $-\lambda$ is therefore the
unique minimum of the linear functional $V\mapsto V_k$ over $\cB_\lambda$'s vertex
set, hence over $\cB_\lambda$ itself. Any $V\in\cB_\lambda$ attaining this minimum
is, by (b), a convex combination of vertices; since every vertex other than
$\{\lambda(e_i-e_k)\}_{i\ne k}$ has $V_k\in\{0,+\lambda\}>-\lambda$, attaining the
minimum $-\lambda$ forces all weight onto $\{\lambda(e_i-e_k)\}_{i\ne k}$. Hence
$\{V\in\cB_\lambda:V_k=-\lambda\}=\mathrm{conv}\{\lambda(e_i-e_k):i\ne
k\}=\lambda\cdot\mathrm{conv}\{e_i-e_k:i\ne k\}$, the image under $u\mapsto u-e_k$
of the standard $(K-2)$-simplex on the $K-1$ coordinates $\ne k$: a
$(K-2)$-simplex.
\end{proof}

\begin{proof}[Proof of Proposition~\ref{prop:exit-interior}]

\emph{Step 1 (rescaling, and which cells' constraints bind).} Write $W=V/\lambda$,
so $W$ is uniform on $\hat F_\lambda:=\hat\cB\cap\{W:W_k\ge-p_0(k)/\lambda\ \forall
k\}$, where $\hat\cB:=\{W\in H:\|W\|_1\le2\}$ is the $\lambda=1$ body of
Lemma~\ref{lem:root-polytope} and does not depend on $\lambda$. Since
$\Sigma(\lambda)=\mathrm{Cov}(V)=\lambda^2\hat\Sigma(\lambda)$ with
$\hat\Sigma(\lambda):=\mathrm{Cov}(W)$, and $\Theta_{ab}=\tau_a^\top\Sigma\tau_b$
enters $\Theta$ only through the scale-invariant ratio
$\Theta_{SY}/\sqrt{\Theta_{SS}\Theta_{YY}}$, $\Theta(\PP_0,\lambda)$ depends on
$\lambda$ only through the \emph{shape} $\hat\Sigma(\lambda)$, not through
$\lambda^2$ itself. For cell $k$, the box constraint on $\hat\cB$ is vacuous
(since $\hat\cB$'s own extent already respects $W_k\ge-1$, by
Lemma~\ref{lem:root-polytope}(a) applied with $S=[K]\setminus\{k\}$, which gives
$\sum_{j\ne k}W_j\le1$, i.e.\ $-W_k\le1$) exactly when $p_0(k)/\lambda\ge1$, i.e.\
$\lambda\le p_0(k)$: this recovers the interior regime's per-cell threshold. At
$\lambda=p_{\min}+\epsilon$ for $\epsilon>0$ small, \emph{every} $k\in\cK^*$
(sharing $p_0(k)=p_{\min}$) has its constraint active at the same depth
$\delta:=1-p_{\min}/\lambda=1-p_{\min}/(p_{\min}+\epsilon)=\epsilon/p_{\min}+O(\epsilon^2)$,
while every $k\notin\cK^*$ has $p_0(k)>p_{\min}$ and so remains strictly interior
($\lambda<p_0(k)$ for $\epsilon$ small enough) with no constraint yet binding.
Fix one $k^*\in\cK^*$ for Steps 2--5 below; Step 6 aggregates over all of $\cK^*$.

\emph{Step 2 (removed volume, to leading order in $\delta$).} By
Lemma~\ref{lem:root-polytope}(c), $\{W\in\hat\cB:W_{k^*}=-1\}$ is the
$(K-2)$-simplex $\mathrm{conv}\{e_i-e_{k^*}:i\ne k^*\}$, a facet of $\hat\cB$ with
positive $(K-2)$-dimensional relative volume. Removing
$R_\delta:=\{W\in\hat\cB:W_{k^*}<-1+\delta\}$ removes a slab adjacent to this
facet, of thickness $\delta$ in the functional $W_{k^*}$ (equivalently, Euclidean
thickness $\delta/\sqrt{1-1/K}$ in the direction $P_He_{k^*}$ normal to the facet
within $H$); writing $A(s):=\mathrm{vol}_{K-2}(\{W\in\hat\cB:W_{k^*}=-1+s\})$ for the
cross-sectional area at level $s$, the coarea formula gives
$\mathrm{vol}(R_\delta)=\int_0^\delta A(s)\,ds$, and since $\hat\cB$ is a polytope,
$A(\cdot)$ is piecewise polynomial and in particular continuous at $s=0$ with
$A(0)>0$ (the facet's own volume); hence
$\rho:=\mathrm{vol}(R_\delta)/\mathrm{vol}(\hat\cB)=\kappa_{\mathrm{vol}}\delta+o(\delta)$
for the constant $\kappa_{\mathrm{vol}}=\kappa_{\mathrm{vol},K}:=A(0)/\mathrm{vol}(\hat\cB)>0$,
which by the $S_K$-symmetry argument of Step 6 does not depend on which cell
$k^*$ is removed.

\emph{Step 3 (second moment on the removed slab).} As $\delta\to0$, the removed
slab's conditional law converges to $\mathrm{Unif}$ on the limiting facet. Writing
$W=u-e_{k^*}$ for $u$ supported on the $K-1$ coordinates $\ne k^*$ summing to $1$,
uniform on the facet means $u$ restricted to those coordinates is
$\mathrm{Dir}(\mathbf1_{K-1})$, with the standard symmetric-Dirichlet second moment
$\E[uu^\top]=\frac1{K(K-1)}(I_{-k^*}+\mathbf1'\mathbf1'^\top)$, where
$I_{-k^*}:=I_K-e_{k^*}e_{k^*}^\top$ and $\mathbf1':=\mathbf1-e_{k^*}$ (both
restricted to, and zero outside, the $K-1$ coordinates $\ne k^*$). Since
$\E[u]=\frac1{K-1}\mathbf1'$,
\begin{equation}\label{eq:Mkstar}
    M_{k^*}:=\lim_{\delta\to0}\E[WW^\top\mid W\in R_\delta]
    = \frac{I_{-k^*}+\mathbf1'\mathbf1'^\top}{K(K-1)}
    -\frac{\mathbf1'e_{k^*}^\top+e_{k^*}\mathbf1'^\top}{K-1}+e_{k^*}e_{k^*}^\top,
\end{equation}
which satisfies $M_{k^*}\mathbf1=0$: $W\in H$ always, so $W^\top\mathbf1=0$
identically, and hence $\E[WW^\top]\mathbf1=\E[W(W^\top\mathbf1)]=0$ on the slab
as well as on all of $\hat\cB$.

\emph{Step 4 (covariance update).} Let $Z\sim\mathrm{Unif}(\hat\cB)$. Since
$V\mapsto-V$ maps $H$ to itself and preserves $\|\cdot\|_1$, it maps $\hat\cB$ to
itself, so $Z\stackrel d=-Z$ and $\E[Z]=0$; by Proposition~\ref{prop:interior}'s
proof, applied directly to $\hat\cB$ (which is $S_K$-invariant for every
$\lambda$, independently of any interior-regime condition, unlike $F_\lambda$
itself), $\Sigma_0:=\mathrm{Cov}(Z)=\E[ZZ^\top]=a_0P_H$ for a fixed constant
$a_0>0$ depending only on $K$. Conditioning $Z$ to avoid $R_\delta$ (probability
mass $\rho$) gives $\E[\hat WW^\top]=\frac{\Sigma_0-\rho M_R}{1-\rho}$, with
$M_R\to M_{k^*}$ as $\delta\to0$; expanding $(1-\rho)^{-1}=1+\rho+O(\rho^2)$,
\[
    \E[\hat W\hat W^\top] = \Sigma_0+\rho(\Sigma_0-M_{k^*})+o(\rho).
\]
The truncated mean $\E[\hat W]=\frac{-\rho\,\E[Z\mid Z\in
R_\delta]}{1-\rho}=O(\rho)$ (bounded by compactness of $\hat\cB$, though not zero:
the removed slab's own centroid need not vanish even though $\E[Z]=0$), so its
contribution to $\hat\Sigma(\lambda)=\E[\hat W\hat
W^\top]-\E[\hat W]\E[\hat W]^\top$ is $O(\rho^2)=o(\rho)$ and does not affect the
following to the order retained:
\begin{equation}\label{eq:sigma-update}
    \hat\Sigma(\lambda)-\Sigma_0=\rho(\Sigma_0-M_{k^*})+o(\rho)
    =\kappa_{\mathrm{vol}}\delta(\Sigma_0-M_{k^*})+o(\delta)
    =\frac{\kappa_{\mathrm{vol}}}{p_{\min}}(\Sigma_0-M_{k^*})\,\epsilon+o(\epsilon).
\end{equation}

\emph{Step 5 (from $\Sigma$ to $\Theta$: a directional derivative of a ratio).}
Write $\Theta(\Sigma)=N(\Sigma)/\sqrt{D_S(\Sigma)D_Y(\Sigma)}$,
$N=\tau_S^\top\Sigma\tau_Y$, $D_a=\tau_a^\top\Sigma\tau_a$; since $A_S,A_Y>0$
(Assumption~\ref{ass:nondegeneracy}'s finite-$K$ form) and $\Theta(\cdot)$ is
continuously differentiable at $\Sigma_0$, \eqref{eq:sigma-update}'s $o(\epsilon)$
remainder in $\hat\Sigma(\lambda)$ propagates to an $o(\epsilon)$ remainder in
$\Theta(\hat\Sigma(\lambda))$, so it suffices to differentiate along the exact
direction $D:=\Sigma_0-M_{k^*}$: writing $\Sigma_0+tD$, $t:=\kappa_{\mathrm{vol}}\epsilon/p_{\min}$,
$A_N=\tau_S^\top\Sigma_0\tau_Y$, $B_N=\tau_S^\top D\tau_Y$ (similarly
$A_S,B_S,A_Y,B_Y$), the same delta-method computation as in the proof of
Theorem~\ref{thm:main} (differentiating $N(t)/\sqrt{D_S(t)D_Y(t)}$ at $t=0$) gives
\begin{equation}\label{eq:gprime}
    \left.\frac{d}{dt}\Theta(\Sigma_0+tD)\right|_{t=0}
    =\frac{B_N}{\sqrt{A_SA_Y}}-\frac{r_K}2\left(\frac{B_S}{A_S}+\frac{B_Y}{A_Y}\right),
    \qquad r_K=\frac{A_N}{\sqrt{A_SA_Y}}=\Corr_K(\tau_S,\tau_Y).
\end{equation}
Because $\Sigma_0\mathbf1=0=M_{k^*}\mathbf1$ (Step 3), $A_N,B_N,A_S,B_S,A_Y,B_Y$
(all six) are unchanged if $\tau_a$ is replaced by its centering
$c_a=\tau_a-\bar\tau_a\mathbf1$; using $P_Hc_a=c_a$, $A_S=a_0\|c_S\|^2$,
$A_N=a_0\langle c_S,c_Y\rangle$, and, from \eqref{eq:Mkstar} together with
$\sum_kc_a(k)=0$ (so $c_a^\top\mathbf1'=-c_a(k^*)$, $c_a^\top e_{k^*}=c_a(k^*)$),
\begin{equation}\label{eq:Mkstar-quad}
    c_S^\top M_{k^*}c_Y=\frac{\langle c_S,c_Y\rangle}{K(K-1)}+\frac{K+1}{K-1}\,c_S(k^*)c_Y(k^*),
\end{equation}
(and likewise for $(S,S)$, $(Y,Y)$). Hence
$B_N=A_N-c_S^\top M_{k^*}c_Y=\beta A_N-\frac{K+1}{K-1}c_S(k^*)c_Y(k^*)$ for the
\emph{same} scalar $\beta=1-\frac1{a_0K(K-1)}$ appearing in
$B_S=\beta A_S-\frac{K+1}{K-1}c_S(k^*)^2$ and the analogous $B_Y$. Substituting into
\eqref{eq:gprime}, the $\beta$-proportional terms cancel exactly
($\beta r_K-\frac{r_K}2(\beta+\beta)=0$, the same degree-zero-homogeneity
cancellation by which $a(\lambda)$ disappears from \eqref{eq:interior-reduction}),
leaving
\[
    \left.\frac{d}{dt}\Theta(\Sigma_0+tD)\right|_{t=0}
    =-\frac{K+1}{K-1}\left[\frac{c_S(k^*)c_Y(k^*)}{\sqrt{A_SA_Y}}
    -\frac{r_K}2\left(\frac{c_S(k^*)^2}{A_S}+\frac{c_Y(k^*)^2}{A_Y}\right)\right]
    =\frac{K+1}{a_0(K-1)}\,\varpi(k^*),
\]
using $A_S=a_0\|c_S\|^2$ and $z_a(k^*)=c_a(k^*)/\|c_a\|$ to rewrite the bracket as
\[
    -\Big[z_S(k^*)z_Y(k^*)-\tfrac{r_K}2\big(z_S(k^*)^2+z_Y(k^*)^2\big)\Big]/a_0=\varpi(k^*)/a_0.
\]
The algebraic equivalence of $\varpi(k)$'s two forms
\eqref{eq:exit-slope}--\eqref{eq:exit-psi-alt} is a direct expansion of
$\tfrac{1+r_K}4(z_S-z_Y)^2-\tfrac{1-r_K}4(z_S+z_Y)^2$ into
$\tfrac{r_K}2(z_S^2+z_Y^2)-z_Sz_Y$, collecting the $z_S^2,z_Y^2,z_Sz_Y$
coefficients. Combining with \eqref{eq:sigma-update}
($t=\kappa_{\mathrm{vol}}\epsilon/p_{\min}$) gives \eqref{eq:exit-slope} for cell
$k^*$ alone, with $C_K:=\kappa_{\mathrm{vol}}(K+1)/(a_0(K-1))>0$.

\emph{Step 6 (tied minimizers: exact additivity).} If $|\cK^*|>1$, the removed
regions $R_\delta^{(k)}:=\{W\in\hat\cB:W_k<-1+\delta\}$, $k\in\cK^*$, are exactly
disjoint for $\delta<1/2$: on $\hat\cB$, $\sum_{j\ne k}W_j\le1$ for every $k$
(Lemma~\ref{lem:root-polytope}(a) with $S=[K]\setminus\{k\}$), so in particular
$W_k+W_{k'}\ge-1$ for any $k\ne k'$; but $W\in
R_\delta^{(k)}\cap R_\delta^{(k')}$ would require $W_k,W_{k'}<-1+\delta$, giving
$W_k+W_{k'}<-2+2\delta<-1$ once $\delta<1/2$, a contradiction. Volumes and
second-moment integrals over $\bigcup_{k\in\cK^*}R_\delta^{(k)}$ are therefore
\emph{exactly} additive (no inclusion--exclusion correction is needed), so
$\hat\Sigma(\lambda)-\Sigma_0=\sum_{k\in\cK^*}\rho_k(\Sigma_0-M_k)+o(\delta)$ with
$\rho_k=\kappa_{\mathrm{vol}}\delta+o(\delta)$ for every $k\in\cK^*$: the constant
$\kappa_{\mathrm{vol}}$ does not depend on which cell is removed, since $\hat\cB$
is invariant under every permutation $\Pi$ of $\{1,\dots,K\}$
(Lemma~\ref{lem:root-polytope}), and $\Pi$ (an orthogonal, volume-preserving map of
$H$) carries $R_\delta^{(k)}$ isometrically onto $R_\delta^{(\Pi(k))}$ and
$M_k$ to $\Pi M_{\Pi(k)}\Pi^\top$. Step 5's computation, being linear in
$D=\Sigma_0-M_k$, applies to each summand separately and adds:
$\partial_\lambda^+\Theta|_{p_{\min}}=\frac{C_K}{p_{\min}}\sum_{k\in\cK^*}\varpi(k)$,
\eqref{eq:exit-slope} as stated. Finally, $\sum_{k=1}^K\varpi(k)=0$: since
$\sum_kz_a(k)^2=\|c_a\|^2/\|c_a\|^2=1$ and
$\sum_kz_S(k)z_Y(k)=\langle c_S,c_Y\rangle/(\|c_S\|\|c_Y\|)=r_K$,
\eqref{eq:exit-slope}'s definition of $\varpi$ gives
$\sum_k\varpi(k)=\frac{r_K}2(1+1)-r_K=0$ directly; when $\cK^*=\{1,\dots,K\}$ (the
uniform-$\PP_0$ case), the right-hand side of \eqref{eq:exit-slope} is exactly this
vanishing sum, consistent with Lemma~\ref{lem:uniform-global}.
\end{proof}

\subsection{Proof of Lemma~\ref{lem:cancellation}}\label{app:proof-cancellation}


\begin{proof}
Write $\mathrm{score}_i$ for the bracketed AIPW expression in \eqref{eq:aipw-estimator}
evaluated at observation $i$. Then
\[
    \hat\Delta_S(\Q_m) = \frac1n\sum_i w_i(\Q_m)\,\mathrm{score}_i
    = \frac1n\sum_{k=1}^K \frac{q_{m,k}}{\hat p_0(k)}\sum_{i:X_i=k}\mathrm{score}_i
    = \sum_{k=1}^K q_{m,k}\cdot\frac{n_k}{n\,\hat p_0(k)}\cdot\frac1{n_k}\sum_{i:X_i=k}\mathrm{score}_i.
\]
Since $\hat p_0(k)=n_k/n$ by definition, $n_k/(n\hat p_0(k))=1$ identically, for every
$k$ with $n_k>0$, regardless of the realized values of $n_k$. Hence
$\hat\Delta_S(\Q_m)=\sum_kq_{m,k}\hat\tau_S(k)$ with $\hat\tau_S(k)$ exactly the
displayed per-cell average, an algebraic identity holding for every finite sample --
the empirical center $\hat p_0$ used to build the importance weight cancels exactly
against the per-cell sample-size normalization used to average within that cell, with
no asymptotic approximation involved.
\end{proof}

\subsection{Proof of Theorem~\ref{thm:main}}\label{app:proof-main}

\begin{proof}
By Lemma~\ref{lem:cancellation}, $\hat\Delta_S(\Q_m)=\sum_kq_{m,k}\hat\tau_S(k)$ where
$\hat\tau_S(k)$ is the per-cell cross-fitted AIPW estimator of $\tau_S(k)$, and
$n_k=\sum_i\mathbf1\{X_i=k\}>0$ for every $k$ with $q_{m,k}>0$ (finitely many cells,
each with $n_k=\Theta_P(n)\to\infty$, so this holds with probability tending to $1$).
Standard cross-fitted AIPW theory (verifying the conditions of
\citealp{chernozhukov2018}, Theorem 3.1: Neyman orthogonality via
Lemma~\ref{lem:if-meanzero}, product-rate nuisance convergence via
Assumption~\ref{ass:nuisance}, cross-fitting to eliminate overfitting bias) gives,
uniformly over $k=1,\dots,K$,
\[
    \hat\tau_S(k)-\tau_S(k) = \frac1{n_k}\sum_{i:X_i=k}\mathrm{score}_S(O_i) + o_P(n^{-1/2}),
    \qquad \mathrm{score}_S(O)=\frac{A(S-\mu_1^S(X))}{e(X)}-\frac{(1-A)(S-\mu_0^S(X))}{1-e(X)}+\mu_1^S(X)-\mu_0^S(X),
\]
the per-unit AIPW score restricted to cell $k$. Multiplying by $q_{m,k}$, summing over
$k$, and using $n_k=n\hat p_0(k)$ together with $w_i(\Q_m)=q_{m,k}/\hat p_0(k)$ for
$X_i=k$ (so that $q_{m,k}/n_k = w_i(\Q_m)/n$ for every $i$ in cell $k$):
\begin{align}
    \hat\Delta_S(\Q_m)-\Delta_S(\Q_m)
    &= \sum_{k=1}^K q_{m,k}\big[\hat\tau_S(k)-\tau_S(k)\big] + o_P(n^{-1/2})
    = \sum_{k=1}^K \frac{q_{m,k}}{n_k}\sum_{i:X_i=k}\mathrm{score}_S(O_i) + o_P(n^{-1/2})\notag\\
    &= \frac1n\sum_{k=1}^K\sum_{i:X_i=k} w_i(\Q_m)\,\mathrm{score}_S(O_i) + o_P(n^{-1/2})
    = \frac1n\sum_{i=1}^n w_i(\Q_m)\,\mathrm{score}_S(O_i) + o_P(n^{-1/2}).
    \label{eq:main-proof-expansion}
\end{align}
Writing this as $\frac1n\sum_i\psi_S^{\mathrm{AIPW}}(O_i;\Q_m)+o_P(n^{-1/2})$ with
$\psi_S^{\mathrm{AIPW}}(O;\Q_m):=w(X;\Q_m)\,\mathrm{score}_S(O)-\Delta_S(\Q_m)$
reproduces \eqref{eq:psi-aipw} exactly, since \eqref{eq:main-proof-expansion}'s
left-hand side is $\hat\Delta_S(\Q_m)-\Delta_S(\Q_m)$: the constant $\Delta_S(\Q_m)$
must be subtracted (rather than any other mean-zero-preserving choice) precisely
because it is what makes $\frac1n\sum_iw_i(\Q_m)\mathrm{score}_S(O_i)$, on the
left of \eqref{eq:main-proof-expansion} before subtraction, equal
$\Delta_S(\Q_m)+o_P(n^{-1/2})$ rather than some other constant --- the expansion
\eqref{eq:main-proof-expansion} is the derivation, not an assertion. (Note that the
alternative centering $w(X;\Q_m)\Delta_S(\Q_m)$ would \emph{also} be mean-zero in
expectation, since $\E_{\PP_0}[w]=1$; mean-zero-ness alone does not distinguish the
two choices. What determines the correct centering is \eqref{eq:main-proof-expansion}
itself: it produces the constant $\Delta_S(\Q_m)$, not $w(X;\Q_m)\Delta_S(\Q_m)$, so
using the latter would introduce the additional term $(w(X;\Q_m)-1)\Delta_S(\Q_m)$,
absent from the true expansion, into the reported influence function.) Direct
verification of mean-zero-ness (Lemma~\ref{lem:if-meanzero}):
$\E_{\PP_0}[w(X;\Q_m)\cdot\mathrm{score}_S(O)]=\E_{\PP_0}[w(X;\Q_m)(\mu_1^S(X)-\mu_0^S(X))]=\E_{\PP_0}[w(X;\Q_m)\tau_S(X)]=\Delta_S(\Q_m)$
by \eqref{eq:delta-linear} and the AIPW residual terms' conditional-mean-zero
property, confirming $\E_{\PP_0}[\psi_S^{\mathrm{AIPW}}(O;\Q_m)]=0$. The analogous
statement holds for $\hat\Delta_Y$.

Correlation is a smooth function $\varphi$ of the five across-study moments $m(\mu)$
(Section~\ref{sec:layer2}'s general statement, specialized), Hadamard differentiable
with bounded gradient on $\{\Var_\mu\Delta_S,\Var_\mu\Delta_Y\ge c_0\}$
(Assumption~\ref{ass:nondegeneracy}). Conditional on $\{\Q_1,\dots,\Q_M\}$, the
functional delta method gives
\[
    \sqrt n\{\Thetahat_n(\lambda)-\Theta(\Qhat,\lambda)\}
    = \frac1{\sqrt n}\sum_{i=1}^n\psi_\Theta(O_i;\Qhat) + o_P(1),
\]
with $\psi_\Theta$ as in \eqref{eq:psi-theta-layer1} (chain rule composing $\varphi$'s
gradient with each per-study score, averaged over $m$). Conditional on the geometry,
$O_1,\dots,O_n$ are i.i.d.\ and independent of $\{\Q_m\}$; $\E_{\PP_0}[\psi_\Theta]=0$
(each stacked score is mean-zero, above) and $\E_{\PP_0}[\psi_\Theta^2]<\infty$ under
bounded weights and bounded gradient (Assumptions~\ref{ass:moments},\ref{ass:nondegeneracy}).
The Lindeberg--L\'evy CLT gives part (a).

For part (b): by Lemma~\ref{lem:mcmc-error} applied to the bounded functional
$g=\varphi(m(\cdot))$ of the chain, $\Theta(\Qhat,\lambda)-\Theta(\PP_0,\lambda)=O_P(M^{-1/2})$,
so $\sqrt n\{\Theta(\Qhat,\lambda)-\Theta(\PP_0,\lambda)\}=O_P((n/M)^{1/2})=o_P(1)$
exactly when $n=o(M)$; combined with part (a) via the triangle inequality (Slutsky),
this gives part (b). Part (c) restates the same decomposition without invoking
$n=o(M)$, retaining both error terms explicitly; the two are asymptotically
independent because the first is a function of $O_1,\dots,O_n$ conditional on the
chain, and the second is a function of the chain alone, and the product-of-orders
$O_P(n^{-1/2})\cdot O_P(M^{-1/2})$ is $o_P$ of either term alone.
\end{proof}

\subsection{Proof of Corollary~\ref{cor:rct}}\label{app:proof-rct}

\begin{proof}
Under (i), $\ehat(X)\equiv\pi$ estimated by the weighted arm-1 proportion
$\bar e_1=\frac1n\sum_iw_i(\Q_m)A_i$; since $\E_{\PP_0}[w(X;\Q_m)]=1$ for any valid
$\Q_m$ (a probability measure), $\bar e_1+\bar e_0\to\E[w A]+\E[w(1-A)]=\E[w]=1$, so
$\bar e_0$ is not a separate estimated quantity and no constant factor of $2$ appears
except in the coincidental case $w\equiv1,\pi=1/2$ where $\bar e_1=\bar e_0=1/2$.
Under (ii), $\hat\mu_1^S(X)\equiv\hat m_{S,1}=\sum_iw_i(\Q_m)A_iS_i/\sum_iw_i(\Q_m)A_i$
by definition of the weighted arm mean, so
$\sum_iw_i(\Q_m)A_i(S_i-\hat m_{S,1})=\sum_iw_i(\Q_m)A_iS_i-\hat m_{S,1}\sum_iw_i(\Q_m)A_i=0$
identically. Substituting into \eqref{eq:aipw-estimator}, the first (IPW-residual)
term for arm 1 vanishes identically, likewise for arm 0, leaving
\[
    \hat\Delta_S(\Q_m) = \frac1n\sum_iw_i(\Q_m)[\hat m_{S,1}-\hat m_{S,0}] = \hat m_{S,1}-\hat m_{S,0}
\]
using $\frac1n\sum_iw_i(\Q_m)=1$ (an identity for any valid weight vector, by
Lemma~\ref{lem:cancellation}'s argument with $\tau_S\equiv$const), exactly
\eqref{eq:iw-estimator}: the estimator collapses exactly as claimed.

For the influence function, we derive $\psi_S^{\mathrm{IW}}$ directly from
\eqref{eq:iw-estimator}'s estimating-equation structure, rather than by substitution
into \eqref{eq:psi-aipw} (substitution is unsound here: the naive claim that the
regression term $\hat m_{S,1}-\hat m_{S,0}$ and the centering
$-(\hat m_{S,1}-\hat m_{S,0})$ are ``equal and opposite'' ignores that
\eqref{eq:psi-aipw}'s regression term enters multiplied by $w(X;\Q_m)$ while the
centering does not, and the two cancel only when $w\equiv1$). The estimator
$\hat m_{S,1}$ solves the $Z$-estimating equation
$\frac1n\sum_iw_i(\Q_m)A_i(S_i-\hat m_{S,1})=0$; by standard $Z$-estimator theory, if
$m_0$ solves the corresponding population equation $\E_{\PP_0}[w(X;\Q_m)A(S-m_0)]=0$
(i.e.\ $m_0=\E[wAS]/\E[wA]$, the population weighted arm-1 mean), then
\[
    \sqrt n(\hat m_{S,1}-m_0) = -\Big(\partial_m\E_{\PP_0}[w(X;\Q_m)A(S-m)]\Big|_{m_0}\Big)^{-1}
    \cdot\frac1{\sqrt n}\sum_iw_i(\Q_m)A_i(S_i-m_0) + o_P(1).
\]
Since $\partial_m\E_{\PP_0}[w(X;\Q_m)A(S-m)]=-\E_{\PP_0}[w(X;\Q_m)A]=-\bar e_1$, this
gives influence function $w(X;\Q_m)A(S-m_0)/\bar e_1$ for $\hat m_{S,1}$; the
analogous computation for arm $0$ gives $-w(X;\Q_m)(1-A)(S-m_0')/\bar e_0$ for
$-\hat m_{S,0}$ (sign flipped since $\hat\Delta_S=\hat m_{S,1}-\hat m_{S,0}$).
Influence functions are additive under subtraction of asymptotically linear
statistics, so
\[
    \psi_S^{\mathrm{IW}}(O;\Q_m) = \frac{w(X;\Q_m)A(S-m_{S,1})}{\bar e_1}
        - \frac{w(X;\Q_m)(1-A)(S-m_{S,0})}{\bar e_0},
\]
exactly \eqref{eq:psi-iw}. This $Z$-estimator argument makes no use of any
cancellation between $w$-weighted and unweighted terms, and is valid for any
$w(X;\Q_m)$, not only $w\equiv1$.
\end{proof}

\subsection{Proof of Lemma~\ref{lem:bilinear-eif}}\label{app:proof-bilinear}

\begin{proof}
Let $\{\PP_t\}$ be a regular parametric submodel through $\PP_0$ at $t=0$ with score
$g(O)$, and let $\tau_{a,t},\tau_{b,t}$ denote the CATEs under $\PP_t$ (holding $C$ and
$d\PP_0$ fixed, per Convention~\ref{conv:geometry}; see
Remark~\ref{rem:kernel-dependence}). Pathwise differentiability of the CATEs
(Lemma~\ref{lem:cate-linearity}) gives, for each fixed $x$,
$\frac{d}{dt}\tau_{a,t}(x)\big|_{t=0}=\E_{\PP_0}[\IF_{\tau_a}(O;x)g(O)]=:\dot\tau_a(x)$,
and similarly for $\tau_b$. Write
$\Theta_{ab}(t)=\iint C(x,x')\tau_{a,t}(x)\tau_{b,t}(x')\,d\PP_0(x)d\PP_0(x')$
(covariate law and $C$ fixed). Differentiating under the integral (Leibniz; justified
because $C$ is bounded and $\tau_{a,t},\tau_{b,t}$ have square-integrable
derivatives, with domination by
$C_w^2(|\dot\tau_a||\tau_b|+|\tau_a||\dot\tau_b|)\in L^1$),
\begin{align*}
    \frac{d}{dt}\Theta_{ab}(t)\Big|_{t=0}
    &= \iint C(x,x')\big[\dot\tau_a(x)\tau_b(x')+\tau_a(x)\dot\tau_b(x')\big]\,d\PP_0(x)d\PP_0(x')\\
    &= \int\dot\tau_a(x)h_b(x)\,d\PP_0(x) + \int\dot\tau_b(x')h_a(x')\,d\PP_0(x'),
\end{align*}
using Fubini to collapse each double integral and defining $h_b,h_a$ as in
\eqref{eq:h-directions}. Substituting $\dot\tau_a(x)=\E_{\PP_0}[\IF_{\tau_a}(O;x)g(O)]$
and interchanging the (finite-measure, bounded-integrand) $x$- and $O$-integrals by
Fubini,
\[
    \frac{d}{dt}\Theta_{ab}(t)\Big|_{t=0}
    = \E_{\PP_0}\!\Big[\Big(\int h_b(x)\IF_{\tau_a}(O;x)\,d\PP_0(x)
        + \int h_a(x')\IF_{\tau_b}(O;x')\,d\PP_0(x')\Big)g(O)\Big].
\]
This exhibits the pathwise derivative as $\E_{\PP_0}[\chi_{ab}(O)g(O)]$ for every
score $g$, so by the Riesz representation of pathwise derivatives, $\chi_{ab}$ (the
bracketed term) is the influence function of $\Theta_{ab}$; it is mean-zero since
$\E_{\PP_0}[\IF_{\tau_a}(O;x)]=0$ for all $x$. Efficiency holds because, in the
nonparametric model, $\chi_{ab}$ is a $\PP_0$-integral of the efficient CATE influence
functions and so lies in the (full) tangent space, making it the canonical gradient.
Setting $a=b$ gives $h_a=h_b$, $\IF_{\tau_a}=\IF_{\tau_b}$, and
$\chi_{aa}=2\int h_a\IF_{\tau_a}\,d\PP_0$. On finite $X$,
$h_b(k)=\sum_{k'}C(k,k')\tau_b(k')p_0(k')=(1/p_0(k))\sum_{k'}\Sigma_{kk'}\tau_b(k')=(\Sigma\tau_b)_k/p_0(k)$
by \eqref{eq:sigma-def}; substituting into $\int h_b\IF_{\tau_a}\,d\PP_0=\sum_kp_0(k)h_b(k)\IF_{\tau_a}(\cdot;k)$
gives $\sum_k(\Sigma\tau_b)_k\IF_{\tau_a}(\cdot;k)=(\Sigma\tau_b)^\top\mathbf{IF}_{\tau_a}$,
reproducing \eqref{eq:chi-ab-discrete}.
\end{proof}

\begin{remark}[Scope: kernel and covariate-law dependence held fixed]\label{rem:kernel-dependence}
Lemma~\ref{lem:bilinear-eif} differentiates only through $(\tau_a,\tau_b)$, holding
$C$ and $d\PP_0$ fixed, matching Theorem~\ref{thm:main}'s conditioning on the sampled
geometry (Convention~\ref{conv:geometry}). An unconditional EIF, which would
additionally differentiate $C$'s and the covariate law's dependence on $\PP_0$ (the
same phenomenon as Remark~\ref{rem:geometry-scope}'s ball-recentering issue), is
deferred to future work.
\end{remark}

\begin{lemma}[Directional representer identity, discrete case]\label{lem:recovery-chi}
On finite $\cX=\{1,\dots,K\}$, $h_b(k)=(\Sigma\tau_b)_k/p_0(k)$ and
$\chi_{ab}(O)=(\Sigma\tau_b)^\top\mathbf{IF}_{\tau_a}(O)+(\Sigma\tau_a)^\top\mathbf{IF}_{\tau_b}(O)$.
\end{lemma}
\begin{proof}
Shown in the proof of Lemma~\ref{lem:bilinear-eif} above, discrete-case paragraph.
\end{proof}

\subsection{Proof of Lemma~\ref{lem:onestep}}\label{app:proof-onestep}

\begin{proof}
By Lemma~\ref{lem:bilinear-eif}, $\chi_{ab}$ is the canonical gradient of $\Theta_{ab}$
in $\beta=(\tau_a,\tau_b)$; adding it to the plug-in cancels the first-order plug-in
bias by the defining property of the one-step construction. Concretely, the
directional derivative of the plug-in at $\beta$ along $(\delta_a,\delta_b)$ is
$\int h_b\delta_a\,d\PP_0+\int h_a\delta_b\,d\PP_0$ (by the product rule computation in
the proof of Lemma~\ref{lem:bilinear-eif}), while $\E_{\PP_0}[\chi_{ab}]$ has
directional derivative equal to its negative; the two cancel, giving orthogonality in
$\beta$. Because $h_b,h_a$ are functions of $x$ that do not depend on the outcome-model
nuisances $\eta=(e,\mu_0,\mu_1)$ (only on $C$ and $\tau$), the ``orthogonality survives
reweighting'' argument of Theorem~\ref{thm:main}'s proof (Appendix~\ref{app:proof-main},
via Lemma~\ref{lem:if-meanzero}) applies verbatim with $h_b$ (resp.\ $h_a$) in the role
of the importance weight $w(X;\Q)$, giving orthogonality in $\eta$ as well.
\end{proof}

\subsection{Proof of Lemma~\ref{lem:elbow-remainder}}\label{app:proof-elbow}

\begin{proof}
The von Mises expansion of \eqref{eq:onestep} has exact second-order remainder
\eqref{eq:remainder}: the one-step correction removes the first-order term
(Lemma~\ref{lem:onestep}), and $\Theta_{ab}$ is exactly bilinear in $\beta$, so the
Taylor expansion terminates at the product term with no higher-order contribution.
Writing $\delta_a=\hat\tau_a-\tau_a$, $\delta_b=\hat\tau_b-\tau_b$,
$R_n=\langle\delta_a,\hat C\delta_b\rangle_{L^2(\PP_0)}$ with
$(\hat C\delta_b)(x)=\int\hat C(x,x')\delta_b(x')\,d\PP_0(x')$; Cauchy--Schwarz and the
operator-norm bound $\|\hat C\|_{op}\le C_w^2$ give
$|R_n|\le\|\hat C\|_{op}\|\delta_a\|\|\delta_b\|$, \eqref{eq:CS-remainder}. Under
H\"older/Sobolev smoothness $s_S,s_Y$ for $\tau_S,\tau_Y$ on $\cX\subseteq\R^d$, the
achievable mean-squared-error-optimal nonparametric rate for $\|\delta_a\|_{L^2}$ is
$n^{-s_a/(2s_a+d)}$, so the product $\|\delta_a\|\|\delta_b\|$ is $o_P(n^{-1/2})$
exactly when $s_a/(2s_a+d)+s_b/(2s_b+d)>1/2$. For $a=b$ (the diagonal moments
$\Theta_{SS},\Theta_{YY}$) this is exactly $s_a>d/2$; for $a\ne b$ it is
$s_S/(2s_S+d)+s_Y/(2s_Y+d)>1/2$, which reduces to the arithmetic $s_S+s_Y>d$ only
under balanced smoothness $s_S=s_Y$. The combined condition $\min(s_S,s_Y)>d/2$
\eqref{eq:elbow-naive} implies both --- the naive sufficient condition of
Remark~\ref{rem:elbow-gap}, not Assumption~\ref{ass:elbow}'s stated
$\min(s_S,s_Y)>d/4$; see that remark for what closing the remaining gap requires.
\end{proof}

\subsection{Proof of Theorem~\ref{thm:general}}\label{app:proof-general}

\begin{proof}
\emph{Step 1 (moment expansions).} For the linear moments,
Lemma~\ref{lem:cate-linearity} integrated against $\bar a$ gives
$\widehat{\E_\mu\Delta_a}-\E_\mu\Delta_a=\frac1n\sum_i\psi_{\E_\mu\Delta_a}(O_i)+o_P(n^{-1/2})$.
For the bilinear moments, Lemma~\ref{lem:elbow-remainder} gives
$\hat\Theta_{ab}-\Theta_{ab}=\frac1n\sum_i\chi_{ab}(O_i)+o_P(n^{-1/2})$ provided
Assumption~\ref{ass:elbow} holds (its only use). Stacking,
$\hat m-m=\frac1n\sum_i\psi_m(O_i)+o_P(n^{-1/2})$.

\emph{Step 2 (delta method).} $\Psi$ is Hadamard differentiable with bounded gradient
on $\{\Var_\mu\Delta_S,\Var_\mu\Delta_Y\ge c_0\}$ (Assumption~\ref{ass:nondegeneracy}),
so
$\sqrt n(\hat\Theta-\Theta(\Qhat))=\nabla\Psi(m)^\top\sqrt n(\hat m-m)+o_P(1)
=\frac1{\sqrt n}\sum_i\psi_\Theta(O_i;\Qhat)+o_P(1)$ by Lemma~\ref{lem:moment-eif}.
For the ratio-type $\Psi$ (correlation), the second-order delta-method remainder is
$O_P(\|\hat m-m\|^2)=O_P(n^{-1})=o_P(n^{-1/2})$ since $\nabla^2\Psi$ is bounded on the
non-degenerate region.

\emph{Step 3 (CLT, conditional on the geometry).} Conditional on the geometry,
$O_1,\dots,O_n$ are i.i.d.\ and independent of it; $\E_{\PP_0}[\psi_\Theta]=0$ (each
stacked score is mean-zero, Lemma~\ref{lem:bilinear-eif}) and
$\E_{\PP_0}[\psi_\Theta^2]<\infty$ under bounded weights, bounded $C$, and bounded
gradient. The Lindeberg--L\'evy CLT gives
$\frac1{\sqrt n}\sum_i\psi_\Theta(O_i;\Qhat)\xrightarrow{d}N(0,\E[\psi_\Theta^2])$.

\emph{Step 4 (Monte Carlo term).} $\Theta(\Qhat)-\Theta(\mu)=O_P(M^{-1/2})$ by
Lemma~\ref{lem:mcmc-error} applied to the bounded functional $g=\Psi(m(\cdot))$ of the
chain; scaling by $\sqrt n$ gives $O_P((n/M)^{1/2})$, which is $o_P(1)$ iff $n=o(M)$.
Combining Steps 1--4 gives \eqref{eq:thm-general}.
\end{proof}

\subsection{Proof of Corollary~\ref{cor:recovery} and Corollary~\ref{cor:onestep-match}}\label{app:proof-recovery}

\begin{proof}[Proof of Corollary~\ref{cor:recovery}]
On finite $X$, Assumption~\ref{ass:elbow} (A5) holds automatically
(Remark~\ref{rem:A5-conservative}) and $\Theta_{ab}=\tau_a^\top\Sigma\tau_b$ by
\eqref{eq:discrete-quadform}. Lemma~\ref{lem:recovery-chi} gives
$\chi_{ab}(O)=(\Sigma\tau_b)^\top\mathbf{IF}_{\tau_a}(O)+(\Sigma\tau_a)^\top\mathbf{IF}_{\tau_b}(O)$
exactly. The per-cell AIPW score is $\IF_{\tau_a}(\cdot;k)$ of Lemma~\ref{lem:cate-linearity}'s
discrete specialization, and $\varphi$'s gradient (correlation's chain-rule
composition through the five moments) is the same object used in
\eqref{eq:psi-theta-layer1}; substituting into \eqref{eq:psi-theta-general}
reproduces $\psi_\Theta(O;\Qhat)$ of Theorem~\ref{thm:main} term-for-term. All
remainders (Lemmas~\ref{lem:cate-linearity},\ref{lem:elbow-remainder}) specialize to
those of Theorem~\ref{thm:main}'s proof, so the two limit laws coincide.
\end{proof}

\begin{lemma}[Plug-in bias of a bilinear functional]\label{lem:plugin-bias}
Suppose, in addition to Lemma~\ref{lem:cate-linearity}'s asymptotic linearity, that
$\E_{\PP_0}[\delta_a(k)]=O(n^{-1})$ for each $k$ (e.g.\ because $\hat\tau_a(k)$ uses
known nuisances, so $\delta_a(k)=\frac1{n_k}\sum_{i:X_i=k}\IF_{\tau_a}(O_i;k)$ exactly,
with $\E_{\PP_0}[\IF_{\tau_a}(O;k)]=0$ exactly and no remainder; for cross-fitted
nuisances this requires the remainder $R_a(k)$ of Lemma~\ref{lem:cate-linearity} to
satisfy the stronger condition $\E_{\PP_0}[R_a(k)]=O(n^{-1})$, not merely
$\|R_a\|_{L^2}=o_P(n^{-1/4})$, which is not implied by
Assumption~\ref{ass:nuisance} alone and is flagged as an additional working
condition here). Then
$\E[\hat\Theta_{ab}^{\mathrm{plugin}}-\Theta_{ab}]=\tr(\Sigma V)/n+O(n^{-3/2})$, with
$V_{kk'}=\E_{\PP_0}[\IF_{\tau_a}(O;k)\IF_{\tau_b}(O;k')]$.
\end{lemma}
\begin{proof}
Write $\hat\tau_a=\tau_a+\delta_a$; since $\Theta_{ab}(\tau)=\tau_a^\top\Sigma\tau_b$
is exactly bilinear, $\hat\Theta_{ab}^{\mathrm{plugin}}-\Theta_{ab}=\tau_a^\top\Sigma\delta_b+\delta_a^\top\Sigma\tau_b+\delta_a^\top\Sigma\delta_b$
exactly, no higher-order terms. Taking expectations, the first two terms are
$\tau_a^\top\Sigma\,\E[\delta_b]$ and $\E[\delta_a]^\top\Sigma\tau_b$, each $O(n^{-1})$
under the stated first-moment condition on $\delta_a,\delta_b$ (not merely $o(n^{-1/2})$,
which would be too weak to be dominated by the $O(n^{-1})$ term below) --- these
$O(n^{-1})$ contributions are absorbed into the lemma's $O(n^{-3/2})$ remainder only
if $\E[\delta_a]=O(n^{-1})$ exactly zero-mean-to-that-order; under the oracle-nuisance
case they vanish identically since $\E[\delta_a(k)]=0$ exactly. The remaining term
satisfies
$\E[\delta_a^\top\Sigma\delta_b]=\sum_{k,k'}\Sigma_{kk'}\E[\delta_a(k)\delta_b(k')]=\sum_{k,k'}\Sigma_{kk'}\cdot\frac1nV_{kk'}+o(n^{-1})=\tr(\Sigma V)/n+O(n^{-3/2})$,
using $\Cov(\delta_a(k),\delta_b(k'))=\frac1n\E[\IF_{\tau_a}(O;k)\IF_{\tau_b}(O;k')]+o(n^{-1})$
from Lemma~\ref{lem:cate-linearity}'s asymptotic linearity and $\Sigma$ symmetric.
\end{proof}

\begin{proof}[Proof of Corollary~\ref{cor:onestep-match} (structural part only; see Remark~\ref{rem:onestep-scope})]
Lemma~\ref{lem:plugin-bias} establishes the plug-in bias $\tr(\Sigma V)/n$ directly.
The one-step estimator \eqref{eq:onestep}'s correction term
$\frac1n\sum_i\hat\chi_{ab}(O_i)$, by Lemma~\ref{lem:onestep}, is constructed
precisely to be the canonical gradient that Neyman-orthogonalizes the plug-in against
first-order perturbations in $(\tau_a,\tau_b)$; since the plug-in's leading bias
$\tr(\Sigma V)/n$ arises exactly from the second-order (product-of-errors) term of the
same bilinear expansion used in Lemma~\ref{lem:plugin-bias}'s proof, and the one-step
correction's defining property (Lemma~\ref{lem:onestep}) is to cancel first-order and
leave only this same second-order remainder (Lemma~\ref{lem:elbow-remainder}), both
corrections target the identical bias term. This establishes the structural
correspondence claimed in Corollary~\ref{cor:onestep-match}; it does not by itself
verify that the package's coded $\tr(\hat\Sigma\hat V)/n$ correction and
\eqref{eq:onestep}'s $n$-vs-$n_k$ renormalized correction agree exactly in finite
samples, which is the point left open in Remark~\ref{rem:onestep-scope}.
\end{proof}

% ===========================================================================
\section{Counterexamples}\label{app:counterex}
% ===========================================================================

\begin{example}[Homogeneous CATE: $\Theta$ undefined]\label{ex:homogeneous}
If $\tau_S(x)$ is constant in $x$ (no effect modification by $X$), then
$\Delta_S(\Q)=\bar\tau_S$ for every $\Q$, so $\Var_\mu\{\Delta_S(\Q)\}=0$ and $\Theta$
is $0/0$, violating Assumption~\ref{ass:nondegeneracy}. This is not pathological: it
is the common scientific case of no heterogeneity in the surrogate's effect and
deserves an explicit diagnostic (report $\Var_\mu\{\Delta_S(\Q)\}$ alongside
$\hat\Theta_n$), not silent failure.
\end{example}

\begin{example}[High PTE, zero $\Theta$]\label{ex:pte-high-theta-zero}
Take $K=3$ equally-weighted cells and CATEs $\tau_S=(1,0,-1)$, $\tau_Y=(1,-2,1)$ (both
non-constant, so neither $\Corr_K$ nor PTE is degenerate). Direct computation gives
$\bar\tau_S=\bar\tau_Y=0$ and $\frac13\sum_k\tau_S(k)\tau_Y(k)=\frac13(1\cdot1+0\cdot(-2)+(-1)\cdot1)=0$,
so $\Corr_K(\tau_S,\tau_Y)=0$ exactly (Proposition~\ref{prop:interior}, interior
regime), with both marginal variances strictly positive
($\frac13\sum_k\tau_S(k)^2=\frac23$, $\frac13\sum_k\tau_Y(k)^2=2$): the correlation is
well-defined, not $0/0$. Meanwhile PTE at $\PP_0$ is a property of a single
(population-weighted) point on the two CATE curves --- the overall mediated share of
the \emph{average} effect --- and can be made close to its maximum by choosing the
direct-effect coefficients so that $S$ carries most of $\PP_0$'s average outcome
effect, independent of how $\tau_S,\tau_Y$ co-vary across strata. A surrogate that
mediates most of the effect in the observed study can therefore be estimated to have
zero across-study transportability: PTE summarizes the effect's decomposition at one
population; $\Theta$ summarizes how the two effects move together across populations,
which depends on the whole shape of $\tau_S,\tau_Y$, not their average level. This
confirms $\Theta$ and PTE answer different questions (Section~\ref{sec:discussion}).
\end{example}

\begin{example}[Binding nonnegativity restores $\lambda$- and $\PP_0$-dependence]\label{ex:binding}
The four canonical DGPs (main text, Table~1) all have $\min_kp_0(k)=0.05$ and use
$\lambda=0.3\gg0.05=\min_kp_0(k)$, so Definition~\ref{def:interior}'s sufficient
condition fails; direct computation (verified numerically against the DGP
specifications) confirms the reported true $\Theta$ values differ from both the
unweighted reduction \eqref{eq:interior-reduction} and the $p_0$-weighted covariance
correlation by amounts well outside numerical precision (e.g.\ DGP1:
$\Theta=0.6907$ vs.\ unweighted $0.6455$ vs.\ $p_0$-weighted $0.6411$), confirming the
general, boundary-aware machinery of Theorem~\ref{thm:main} is doing genuine work in
the validated regime, not merely recovering the closed form.
\end{example}

\begin{example}[Below the elbow: no $\sqrt n$-normal estimator]\label{ex:below-elbow}
Take $\cX=[0,1]^d$ with $d=10$ and $\tau_S,\tau_Y$ each Lipschitz ($s_S=s_Y=1$) but no
smoother. Then $\min(s_S,s_Y)=1<d/4=2.5$: Assumption~\ref{ass:elbow} fails (the
balanced case, where diagonal and off-diagonal boundaries coincide), the minimax
rate for $\Theta_{SS}$ (or $\Theta_{YY}$) is $n^{-4/14}$, strictly slower than
$n^{-1/2}$, and no $\sqrt n$-consistent, asymptotically normal estimator of $\Theta$
exists regardless of estimator choice \citep{BickelRitov1988,Robins2008HOIF};
discretizing $X$ into a coarser grid (reducing effective $d$, or moving toward the
finite-rank endpoint of Section~\ref{sec:rank}) is the practical remedy, at the cost
of the information lost to discretization.

The diagonal condition binds independently of the off-diagonal one: take $d=8$,
$s_S=0.9$, $s_Y=3.2$. Then $s_S+s_Y=4.1>4=d/2$ (the off-diagonal boundary holds),
but $s_S=0.9<2=d/4$ (the diagonal boundary for $\Theta_{SS}$ fails), so
$\min(s_S,s_Y)=0.9<d/4$: Assumption~\ref{ass:elbow} still fails, driven entirely by
$\Theta_{SS}$'s minimax rate $n^{-4(0.9)/(4(0.9)+8)}=n^{-3.6/11.6}$, even though the
cross-term $\Theta_{SY}$ alone would already be $\sqrt n$-estimable at these
smoothness values. This is why the combined condition is $\min(s_S,s_Y)>d/4$ and not
the (weaker, insufficient) off-diagonal condition $s_S+s_Y>d/2$ alone.
\end{example}

\begin{thebibliography}{9}
\bibitem[Bickel \& Ritov(1988)]{BickelRitov1988} Bickel, P. J. and Ritov, Y. (1988).
  Estimating integrated squared density derivatives: sharp best order of convergence
  estimates. \emph{Sankhy\={a} A} 50, 381--393.
\bibitem[Birg\'e \& Massart(1995)]{BirgeMassart1995} Birg\'e, L. and Massart, P.
  (1995). Estimation of integral functionals of a density. \emph{Annals of
  Statistics} 23, 11--29.
\bibitem[Chernozhukov et al.(2018)]{chernozhukov2018} Chernozhukov, V., Chetverikov,
  D., Demirer, M., Duflo, E., Hansen, C., Newey, W. and Robins, J. (2018).
  Double/debiased machine learning for treatment and structural parameters.
  \emph{The Econometrics Journal} 21, C1--C68.
\bibitem[Devlin et al.(1975)]{DevlinGnanadesikanKettenring1975} Devlin, S. J.,
  Gnanadesikan, R. and Kettenring, J. R. (1975). Robust estimation and outlier
  detection with correlation coefficients. \emph{Biometrika} 62, 531--545.
\bibitem[Kennedy et al.(2024)]{Kennedy2024Minimax} Kennedy, E. H., Balakrishnan, S.,
  Robins, J. M. and Wasserman, L. (2024). Minimax rates for heterogeneous causal
  effect estimation. \emph{Annals of Statistics} 52, 793--816.
\bibitem[Postnikov(2009)]{Postnikov2009} Postnikov, A. (2009). Permutohedra,
  associahedra, and beyond. \emph{International Mathematics Research Notices} 2009,
  1026--1106.
\bibitem[Robins et al.(2008)]{Robins2008HOIF} Robins, J., Li, L., Tchetgen, E. and van
  der Vaart, A. (2008). Higher order influence functions and minimax estimation of
  nonlinear functionals. \emph{Probability and Statistics: Essays in Honor of David
  A. Freedman} 2, 335--421.
\bibitem[Robins et al.(2017)]{Robins2017Minimax} Robins, J. M., Li, L., Mukherjee, R.,
  Tchetgen Tchetgen, E. and van der Vaart, A. (2017). Minimax estimation of a
  functional on a structured high-dimensional model. \emph{Annals of Statistics} 45,
  1951--1987.
\bibitem[Stone(1980)]{Stone1980} Stone, C. J. (1980). Optimal rates of convergence
  for nonparametric estimators. \emph{Annals of Statistics} 8, 1348--1360.
\bibitem[Williams(1991)]{Williams1991} Williams, D. (1991). \emph{Probability with
  Martingales}. Cambridge Mathematical Textbooks. Cambridge University Press.
\end{thebibliography}

\end{document}
